\chapter{Teoria węzłów: od diagramów do rozmaitości czterowymiarowych}
\label{ch:teoria-wezlow}

\noindent\textbf{Związek z wcześniejszym materiałem.}
Przy wprowadzaniu teorii węzłów wykorzystujemy poznaną teorię
grafów planarnych i wielomian Tutte'a. Po podstawach teorii
węzłów następuje nowy rozdział pokazujący bezpośrednio, jak
z grafów Taita wyprowadzić niezmienniki Jonesa i Goeritza.
W dalszych rozdziałach skorzystamy z teorii Bassa--Serre'a,
chirurgii Dehna i rozkładów JSJ.


Teoria węzłów łączy elementarne pytania o zanurzenia okręgu w przestrzeni
trójwymiarowej z fundamentalnymi problemami geometrii i topologii.
Jeden i ten sam węzeł można opisać diagramem na płaszczyźnie, grupą
fundamentalną jego uzupełnienia, wielomianami i grupami homologii,
strukturą geometryczną uzupełnienia lub sposobem dołączenia uchwytu
wymiaru cztery. Celem rozdziału jest przeprowadzenie tych przejść,
z precyzyjnym rozróżnieniem argumentów przeprowadzonych w całości
i twierdzeń wymagających odrębnej, zaawansowanej teorii. W drugim
przypadku formułujemy dokładną tezę, określamy używane narzędzie
i wyjaśniamy zależność od wcześniejszego materiału; szkic nie jest
przedstawiany jako kompletny dowód.

Przyjmujemy kategorię gładką, o ile nie zaznaczono inaczej.
Rozmaitości i sploty są orientowane, jeśli orientacja jest potrzebna.
Sfera $S^3$ jest orientowanym brzegiem kuli $B^4$; utożsamiamy ją także
z jednopunktowym uzwarceniem $\mathbb R^3$. W szczególności
nie odróżniamy węzłów w $\mathbb R^3$ i w $S^3$, o ile ich nośniki
nie przechodzą przez punkt w nieskończoności.

\section{Węzły, sploty i izotopie}
\label{sec:kt-podstawy}

\begin{uwaga}[Konwencja terminologiczna]
W całym skrypcie słowa \emph{zanurzenie} używamy w znaczeniu
angielskiego \emph{embedding}: gładkie odwzorowanie
$f\colon N\to M$ jest zanurzeniem, gdy jest injektywną immersją
i homeomorfizmem $N$ na $f(N)$ z topologią podprzestrzeni.
W szczególności różniczka $\mathrm df_x$ jest injektywna w każdym
$x\in N$. Samo założenie injektywności różniczki definiuje
\emph{immersję} i nie wystarcza do zanurzenia w powyższym
znaczeniu. Jeżeli $N$ jest zwarte, injektywna immersja
$N\to M$ do rozmaitości Hausdorffa jest zanurzeniem.
Odwołując się do kategorii topologicznej, przez zanurzenie
lokalnie płaskie rozumiemy \emph{lokally flat embedding}.
\end{uwaga}


\begin{definicja}[Węzeł i splot]
Niech $M$ będzie gładką rozmaitością trójwymiarową.
\emph{Węzłem} w $M$ nazywamy obraz gładkiego zanurzenia
$k\colon S^1\hookrightarrow M$.
\emph{Splotem o $r$ składowych} nazywamy obraz gładkiego zanurzenia
$k\colon\bigsqcup_{i=1}^{r}S^1\hookrightarrow M$, gdzie $r\geq1$.
Splot $L$ jest \emph{zorientowany}, jeżeli na każdej jego składowej
wybrano orientację. Węzeł \emph{trywialny} $U\subset S^3$ to okrąg
ograniczający gładko zanurzony dysk w $S^3$. Splot trywialny o $r$
składowych to splot będący brzegiem $r$ parami rozłącznych,
gładko zanurzonych dysków w $S^3$.
\end{definicja}

Zastrzeżenie dotyczące dysków jest istotne: każda pojedyncza składowa
splotu może być węzłem trywialnym, mimo że cały splot jest
nietrywialny. Pierwszym przykładem jest splot Hopfa.

\begin{definicja}[Izotopia otoczenia]
Niech $K_0,K_1\subset M$ będą węzłami. \emph{Izotopią otoczenia}
przenoszącą $K_0$ na $K_1$ nazywamy gładkie odwzorowanie
$H\colon M\times[0,1]\to M$ takie, że $H_t:=H(\cdot,t)$
jest dyfeomorfizmem dla każdego $t$, $H_0=\mathrm{id}_M$ i
$H_1(K_0)=K_1$. Węzły nazywamy \emph{równoważnymi}, jeżeli
istnieje taka izotopia. Dla splotów definicja jest analogiczna,
z zachowaniem oznaczeń i orientacji składowych, jeśli są ustalone.
\end{definicja}

Twierdzenie o rozszerzaniu izotopii zanurzeń stwierdza precyzyjniej,
że dla gładkiej rodziny zanurzeń $k_t\colon S^1\hookrightarrow M$,
$t\in[0,1]$, istnieje izotopia otoczenia $H_t$ taka, że
$k_t=H_t\circ k_0$ dla wszystkich $t$. Twierdzenie stosuje się,
gdy źródło zanurzenia jest zwarte, a rodzina jest gładka względem
obu zmiennych. Jest to wynik niebanalny w ogólnej teorii podrozmaitości,
ale bardzo wygodny w praktyce. Wszystkie węzły rozpatrujemy zatem
\emph{z dokładnością do izotopii otoczenia}, nie tylko do homeomorfizmu
ich abstrakcyjnych nośników, które zawsze są okręgami.

\begin{przyklad}[Węzły torusowe]
Niech $p,q\in\mathbb Z$ będą względnie pierwsze. Na standardowym
torusie $T^2\subset S^3$, parametryzowanym przez
$S^1\times S^1$, rozważmy krzywą
$t\mapsto(e^{ipt},e^{iqt})$, gdzie $t\in\mathbb R/2\pi\mathbb Z$.
Jej obraz, po zanurzeniu standardowego torusa w $S^3$, oznaczamy
$T(p,q)$ i nazywamy \emph{węzłem torusowym}. Jeśli $|p|=1$
lub $|q|=1$, jest on trywialny. Węzeł $T(2,3)$ jest trójlistnikiem;
węzeł $T(2,-3)$ jest jego lustrzanym odbiciem, z odpowiednio
dobraną konwencją zanurzenia torusa.
\end{przyklad}

\begin{definicja}[Lustrzane odbicie i odwracalność]
Niech $K\subset S^3$ będzie zorientowanym węzłem.
Przez $\overline K$ oznaczamy jego obraz przy dyfeomorfizmie
$S^3\to S^3$ odwracającym orientację, z orientacją na obrazie
przeniesioną z $K$. Przez $K^{\mathrm{rev}}$ oznaczamy ten sam
węzeł z przeciwną orientacją. Węzeł jest \emph{amfichiralny},
gdy jest równoważny swemu odbiciu lustrzanemu po pominięciu
orientacji; jest \emph{odwracalny}, gdy $K$ i $K^{\mathrm{rev}}$
są równoważne jako węzły zorientowane.
\end{definicja}

\begin{uwaga}
Nie należy utożsamiać lustrzanego odbicia z odwróceniem orientacji.
Przykładowo trójlistnik jest odwracalny, ale nie jest amfichiralny.
Wykazanie drugiej własności wymaga niezmiennika rozróżniającego
jego dwa odbicia, np. sygnatury lub wielomianu Jonesa.
\end{uwaga}


\begin{center}
\begin{tikzpicture}[scale=.9,>=Stealth,line cap=round]
  \begin{scope}[shift={(-2.65,0)}]
    \fill[ktBlue!7] (0,0) circle (1.22);
    \draw[ktBlue,line width=3.5pt,-{Stealth[length=3mm]}] (0,0) circle (1.0);
    \node[text=ktInk,font=\small\bfseries] at (0,-1.55) {węzeł trywialny};
    \node[text=ktBlue,font=\small] at (0,-1.95) {$U$};
  \end{scope}
  \begin{scope}[shift={(2.65,0)}]
    \fill[ktTeal!7] (-.9,0) circle (.77);
    \fill[ktCoral!7] (.9,0) circle (.77);
    \draw[ktTeal,line width=3pt] (-.9,0) circle (.74);
    \draw[ktCoral,line width=3pt] (.9,0) circle (.74);
    \node[text=ktInk,font=\small\bfseries] at (0,-1.55) {splot rozłączny};
    \node[text=ktCoral,font=\small] at (0,-1.95) {$U\sqcup U$};
  \end{scope}
\end{tikzpicture}
\end{center}
\noindent Po prawej przedstawiono dwie rozłączne, niesplecione składowe.
Kolor oznacza przynależność do jednej składowej splotu.

\section{Diagramy i ruchy Reidemeistera}
\label{sec:kt-diagramy}

Wybierzmy standardowe rzutowanie $\pi\colon\mathbb R^3\to\mathbb R^2$.
Po dowolnie małym zaburzeniu zanurzenia możemy założyć, że $\pi|_K$
jest immersją poza skończoną liczbą punktów podwójnych, że wszystkie
przecięcia są poprzeczne oraz że nie występują punkty potrójne.
Zaznaczając przy każdym punkcie podwójnym, które z dwóch włókien
leży wyżej, otrzymujemy \emph{diagram węzła}. Odcinek przebiegający
pod drugim rysujemy z krótką przerwą. Rozważania są identyczne dla
splotów.

\begin{center}
\begin{tikzpicture}[scale=.88,line cap=round,line join=round,>=Stealth]
\begin{scope}[shift={(-3.5,0)}]
  \node[font=\small\bfseries,text=ktInk] at (0,1.55) {pasmo niebieskie nad czerwonym};
  \draw[ktCoral,line width=2.8pt] (-1.0,1.0)--(-.14,.14);
  \draw[ktCoral,line width=2.8pt] (.14,-.14)--(1,-1);
  \draw[ktBlue,line width=3.2pt,-{Stealth[length=3mm]}] (-1,-1)--(1,1);
  \draw[ktCoral,line width=2.8pt,-{Stealth[length=3mm]}] (-1,1)--(-.44,.44);
\end{scope}
\begin{scope}[shift={(3.5,0)}]
  \node[font=\small\bfseries,text=ktInk] at (0,1.55) {pasmo czerwone nad niebieskim};
  \draw[ktBlue,line width=2.8pt] (-1,-1)--(-.14,-.14);
  \draw[ktBlue,line width=2.8pt] (.14,.14)--(1,1);
  \draw[ktCoral,line width=3.2pt,-{Stealth[length=3mm]}] (-1,1)--(1,-1);
  \draw[ktBlue,line width=2.8pt,-{Stealth[length=3mm]}] (-1,-1)--(-.43,-.43);
\end{scope}
\node[font=\footnotesize,text=ktInk] at (0,-1.55)
  {Przerwa na diagramie oznacza pasmo przebiegające pod spodem.};
\end{tikzpicture}
\end{center}

\begin{definicja}[Liczba skrzyżowań]
Niech $D$ będzie diagramem splotu $L$. Oznaczamy przez $c(D)$ liczbę
punktów podwójnych diagramu. \emph{Liczbą skrzyżowań} splotu jest
$c(L):=\min\{c(D):D\text{ jest diagramem }L\}$. Diagram realizujący
minimum nazywamy \emph{minimalnym}.
\end{definicja}

\begin{twierdzenie}[Reidemeister]
Niech $D_0,D_1$ będą dwoma regularnymi diagramami splotów
na zorientowanej płaszczyźnie $\mathbb R^2$, wraz z informacją,
która gałąź przechodzi górą w każdym skrzyżowaniu. Diagramy reprezentują równoważne sploty wtedy
i tylko wtedy, gdy jeden z nich można przekształcić w drugi
skończonym ciągiem izotopii płaszczyzny oraz trzech lokalnych
ruchów Reidemeistera: $\mathrm{R I}$ (dodanie lub usunięcie
pętli z jednym skrzyżowaniem), $\mathrm{R II}$ (dodanie lub
usunięcie dwóch skrzyżowań) i $\mathrm{R III}$ (przesunięcie
jednego pasma nad skrzyżowaniem dwóch pozostałych).
\end{twierdzenie}

\begin{proof}[Uzasadnienie geometryczne; użycie transwersalności]
Każdy z ruchów realizuje się w małej kuli zawierającej wyłącznie
odpowiednie fragmenty zanurzenia. Na przykład w ruchu $\mathrm{R II}$
jeden łuk przesuwa się nad drugim, a w $\mathrm{R III}$ pasmo
przechodzi nad skrzyżowaniem pozostałych dwóch. Przedłużamy
te lokalne izotopie przez identyczność poza kulą. Otrzymujemy
izotopie otoczenia, więc operacje zachowują typ splotu.

Dla kierunku odwrotnego nie wystarcza arbitralne oglądanie
rysunków w kolejnych chwilach. Ustalmy gładką izotopię
$H_t$ przenoszącą $L_0$ na $L_1$. Złożenie
$\pi\circ H_t|_{L_0}$ jest jednoparametrową rodziną
odwzorowań krzywych do płaszczyzny. Stosujemy parametryczną
transwersalność: po zaburzeniu rodziny, nieruszając jej
końców, wszystkie rzuty poza skończoną liczbą chwil mają
wyłącznie poprzeczne punkty podwójne o różnych wysokościach.
W chwili wyjątkowej występuje dokładnie jedna osobliwość
najniższej kodymensji. Są trzy możliwości: styczność kierunku
rzutowania do łuku (powstanie lub zanik pętli), styczność
dwóch gałęzi obrazu (powstanie lub zanik pary punktów
podwójnych) albo przecięcie trzech gałęzi obrazu. Odpowiadają
one ruchom $\mathrm{R I}$, $\mathrm{R II}$ i $\mathrm{R III}$.

Między chwilami wyjątkowymi diagramy zmieniają się wyłącznie
izotopią płaszczyzny. Zwartość przedziału parametru i ogólne
położenie pozwalają wybrać skończoną liczbę chwil wyjątkowych,
a więc skończony ciąg ruchów. Przywołana klasyfikacja
osobliwości rodziny rzutowań jest jedynym krokiem dowodu,
w którym używamy parametrycznej transwersalności.
\end{proof}


\begin{center}
\begin{tikzpicture}[scale=.86,line cap=round,line join=round,>=Stealth]
  \node[text=ktInk,font=\footnotesize\bfseries] at (-3.32,1.62) {R I: pętla};
  \draw[ktBlue,line width=2.8pt] (-5.3,-1)--(-5.3,1);
  \draw[->,ktTeal,line width=2pt] (-4.75,0)--(-3.95,0);
  \draw[ktBlue,line width=2.8pt] (-2.4,-1) .. controls (-2.7,-.2)
    and (-1.75,-.4) .. (-1.55,.18)
    .. controls (-1.4,.83) and (-2.65,1.4) .. (-2.8,.45)
    .. controls (-2.9,-.2) and (-2.4,-.65) .. (-2.4,1);
  \node[text=ktInk,font=\footnotesize\bfseries] at (2.64,1.62) {R II: dwa skrzyżowania};
  \draw[ktCoral,line width=2.8pt] (.3,-1)--(.3,1);
  \draw[ktBlue,line width=2.8pt] (1.0,-1)--(1.0,1);
  \draw[->,ktTeal,line width=2pt] (1.5,0)--(2.32,0);
  \draw[ktCoral,line width=2.8pt] (2.9,-1)
    .. controls (4.6,-.8) and (4.7,-.25) .. (3.9,.05)
    .. controls (3.0,.32) and (2.8,.8) .. (2.9,1);
  \draw[ktBlue,line width=2.8pt] (3.7,-1)--(3.7,1);
  \draw[white,line width=6pt] (3.95,.035)
    .. controls (3.78,.12) and (3.57,.22) .. (3.46,.26);
  \draw[ktCoral,line width=2.8pt] (3.95,.035)
    .. controls (3.78,.12) and (3.57,.22) .. (3.46,.26);
  \node[text=ktInk,font=\footnotesize] at (0,-1.52)
    {Schematy lokalne: izotopia w przestrzeni nie wymaga przecinania pasm.};
\end{tikzpicture}
\end{center}

Twierdzenie Reidemeistera jest fundamentem konstrukcji niezmienników
diagramowych. Wystarczy zdefiniować wielkość dla diagramów, a następnie
sprawdzić jej niezmienniczość przy trzech rodzajach ruchów.

\begin{definicja}[Znak skrzyżowania i skręt]
Niech $D$ będzie zorientowanym diagramem. W punkcie skrzyżowania
$c$ wybierzmy dodatnio zorientowane współrzędne na płaszczyźnie.
Niech $v_{\mathrm{nad}},v_{\mathrm{pod}}$ oznaczają wektory styczne
do zorientowanych pasm, odpowiednio przechodzącego górą i dołem.
Definiujemy $\varepsilon(c):=\operatorname{sgn}\det(v_{\mathrm{pod}},v_{\mathrm{nad}})
\in\{-1,1\}$ oraz \emph{skręt diagramu}
$w(D):=\sum_{c}\varepsilon(c)$.
\end{definicja}

Wielkość $w(D)$ pozostaje niezmienna przy ruchach $\mathrm{R II}$
i $\mathrm{R III}$, lecz przy $\mathrm{R I}$ zmienia się o $\pm1$.
Nie jest więc niezmiennikiem węzła. Będzie natomiast koniecznym
czynnikiem normalizującym wielomian Jonesa.

\begin{przyklad}[Dlaczego sama liczba skrzyżowań nie wystarcza]
Na diagramie okręgu można wykonać dowolnie wiele ruchów $\mathrm{R I}$.
Każdy nowy diagram przedstawia węzeł trywialny, chociaż liczba
skrzyżowań jest coraz większa. Nie oznacza to, że $c(K)$ nie jest
niezmiennikiem: minimum po wszystkich diagramach nie zależy od
wybranego przedstawienia.
\end{przyklad}

\section{Liczba splecenia i dualność Alexandera}
\label{sec:kt-linking}

\begin{definicja}[Liczba splecenia]
Niech $K_1,K_2\subset S^3$ będą rozłącznymi, zorientowanymi węzłami.
Wybieramy zorientowaną, zwartą, gładko zanurzoną powierzchnię
$F\subset S^3$ o brzegu $\partial F=K_1$, poprzeczną do $K_2$.
\emph{Liczbą splecenia} nazywamy algebraiczną liczbę punktów
przecięcia $\operatorname{lk}(K_1,K_2):=F\cdot K_2$.
Znak punktu przecięcia wyznaczamy przez orientacje $F$, $K_2$ i $S^3$.
\end{definicja}

Istnienie powierzchni $F$ zostanie udowodnione konstruktywnie
w części poświęconej algorytmowi Seiferta. Następująca własność
pokazuje, że definicja nie zależy od wyboru $F$.

\begin{stwierdzenie}
Dla dowolnych rozłącznych, zorientowanych węzłów
$K_1,K_2\subset S^3$ liczba $\operatorname{lk}(K_1,K_2)$ jest
dobrze określona, symetryczna względem składników i niezmiennicza
przy izotopii zorientowanego splotu.
\end{stwierdzenie}

\begin{proof}
Jeżeli $F$ i $F'$ są dwoma łańcuchami powierzchniowymi z tym samym
brzegiem $K_1$, to $F-F'$ jest cyklem dwuwymiarowym w $S^3$.
Ponieważ $H_2(S^3;\mathbb Z)=0$, jest brzegiem łańcucha
trójwymiarowego. Algebraiczne przecięcie brzegu łańcucha z zamkniętą
krzywą $K_2$ wynosi zero, a więc $F\cdot K_2=F'\cdot K_2$.
Dla symetrii wybieramy powierzchnie $F_1,F_2$ z brzegami
$K_1,K_2$ i ustawiamy je poprzecznie. Ich przecięcie jest
zorientowaną jednowymiarową rozmaitością z brzegiem, a rachunek
orientacji na końcach tego przecięcia daje
$F_1\cdot K_2=F_2\cdot K_1$. Niezmienniczość przy izotopii wynika z
zachowania algebraicznej liczby przecięć w rodzinie poprzecznej.
\end{proof}

\begin{twierdzenie}[Wzór diagramowy]
Niech $D$ będzie zorientowanym diagramem dwuskładowego splotu
$K_1\cup K_2\subset S^3$. Niech $C_{12}(D)$ oznacza zbiór
skrzyżowań, w których uczestniczą różne składowe. Wówczas
\[
\operatorname{lk}(K_1,K_2)=\frac12\sum_{c\in C_{12}(D)}\varepsilon(c).
\]
W szczególności suma znaków skrzyżowań między dwiema składowymi
jest parzysta.
\end{twierdzenie}

\begin{proof}
Oznaczmy prawą stronę wzoru przez $I(D)$. W ruchu $\mathrm{R I}$
na jednej składowej nie powstaje skrzyżowanie międzyskładowe.
W ruchu $\mathrm{R II}$ dwa ewentualne skrzyżowania różnych
składowych mają przeciwne znaki. Ruch $\mathrm{R III}$ nie
zmienia znaków ani par składowych uczestniczących w
skrzyżowaniach. Zatem $I(D)$ zależy tylko od izotopii splotu.

Zmienienie jednego skrzyżowania między $K_1$ i $K_2$
zmienia $I(D)$ o $1$ albo $-1$. Taka sama zmiana zachodzi
w $F\cdot K_2$: w małej kuli realizującej zmianę
skrzyżowania krzywa $K_2$ przechodzi na przeciwną stronę
powierzchni rozpiętej na $K_1$, przekraczając ją dokładnie
raz z odpowiednim znakiem. Zmiana skrzyżowania wewnątrz
jednej składowej nie wpływa na żadną z obu wielkości.

Dowolny splot dwóch składowych można uczynić rozdzielnym
przez skończony ciąg zmian skrzyżowań: wybieramy najpierw
kolejność składowych, przeprowadzamy jedną z nich zawsze
ponad drugą, a następnie umieszczamy je w rozłącznych kulach.
Dla splotu rozdzielnego mamy $I(D)=0$ i
$\operatorname{lk}(K_1,K_2)=0$. Cofając zmiany skrzyżowań,
wnioskujemy, że obie wielkości są równe dla wyjściowego
splotu. Ponieważ liczba splecenia jest całkowita, suma znaków
międzyskładowych jest parzysta.
\end{proof}

\begin{przyklad}
Dla splotu Hopfa o zgodnych orientacjach dwóch skrzyżowań
otrzymujemy $\operatorname{lk}=1$. Dla splotu rozdzielnego
liczba splecenia wynosi zero. Odwrotna implikacja jest fałszywa:
trójskładowe pierścienie Borromeusza mają zerowe liczby splecenia
każdej pary składowych, ale cały splot jest nierozdzielny.
\end{przyklad}

\begin{twierdzenie}[Homologia uzupełnienia splotu]
Niech $L=K_1\cup\dots\cup K_r\subset S^3$ będzie splotem o $r\geq1$
składowych i niech $X_L:=S^3\setminus\operatorname{int}N(L)$,
gdzie $N(L)$ oznacza sumę rozłącznych, domkniętych otoczeń
tubularnych składowych. Wtedy
$H_1(X_L;\mathbb Z)\cong\mathbb Z^r$.
Izomorfizm wyznaczają klasy dodatnio zorientowanych meridianów
poszczególnych składowych.
\end{twierdzenie}

\begin{proof}
Przestrzeń $X_L$ jest homotopijnie równoważna $S^3\setminus L$.
Z dualności Alexandera otrzymujemy
$\widetilde H_i(S^3\setminus L;\mathbb Z)
\cong\widetilde H^{2-i}(L;\mathbb Z)$.
Dla $i=1$ prawa strona jest $H^1(L;\mathbb Z)\cong\mathbb Z^r$.
Meridian otaczający $K_i$ ma liczbę splecenia jeden z $K_i$
i zero z pozostałymi składowymi. Homomorfizm przypisujący
klasie pętli jej $r$ liczb splecenia identyfikuje więc
meridiany z bazą standardową, kończąc dowód.
\end{proof}

\section{Grupa węzła i prezentacja Wirtingera}
\label{sec:kt-wirtinger}

\begin{definicja}
Dla węzła $K\subset S^3$ \emph{grupą węzła} nazywamy
$G_K:=\pi_1(S^3\setminus K,x_0)$, gdzie $x_0\in S^3\setminus K$.
Zmiana punktu bazowego daje grupę izomorficzną, lecz niekanonicznie.
Dla splotu $L$ analogicznie oznaczamy $G_L:=\pi_1(S^3\setminus L,x_0)$.
\end{definicja}

Rysunek diagramu koduje prezentację tej grupy. Dzielimy go na
\emph{łuki Wirtingera}: maksymalne fragmenty przebiegające
bez przerwania pod innym pasmem. Każdemu takiemu łukowi
przypisujemy generator, reprezentowany przez meridian.

\begin{twierdzenie}[Prezentacja Wirtingera]
Niech $D$ będzie spójnym diagramem węzła $K\subset S^3$,
z co najmniej jednym skrzyżowaniem, a $x_1,\dots,x_n$ niech będą
meridianami odpowiadającymi łukom Wirtingera. W każdym skrzyżowaniu
trzy uczestniczące generatory spełniają relację postaci
$x_j=x_i x_k x_i^{-1}$ albo $x_j=x_i^{-1}x_kx_i$,
w zależności od orientacji i położenia pasm; $x_i$ odpowiada
pasmowi górnemu. Grupa $G_K$ ma prezentację o tych generatorach
i relacjach, przy czym jedna relacja jest konsekwencją pozostałych.
\end{twierdzenie}

\begin{proof}[Konstrukcja van Kampena]
Wybieramy regularny diagram i podnosimy łuki górne na
niewielką dodatnią wysokość względem płaszczyzny rzutowania.
Przestrzeń nad diagramem po usunięciu niewielkich rurek
wokół pasm deformacyjnie retraktuje się na układ, którego
grupa fundamentalna jest wolna na meridianach $x_1,\dots,x_n$.
Dołączając części przestrzeni pod diagramem, w każdym
skrzyżowaniu narzucamy jedną relację. Jeśli $x_i$ jest
meridianem łuku górnego, a $x_k,x_j$ oznaczają meridiany
łuków dolnych, przeciągnięcie małego meridianu spod pasma
górnego jest sprzężeniem, zatem $x_j=x_i^{\pm1}x_kx_i^{\mp1}$.
Twierdzenie van Kampena pokazuje, że nie ma innych relacji.

Żeby wyjaśnić redundancję, wybierzmy sferę obejmującą
cały diagram oraz jeden wspólny punkt bazowy dla meridianów.
Odpowiadające skrzyżowaniom relacje są brzegami
komórek dwuwymiarowych standardowego kompleksu Wirtingera.
Suma tych komórek tworzy zamkniętą powierzchnię po uwzględnieniu
orientacji łuków; stąd jedna relacja jest konsekwencją
pozostałych. Formalnie można usunąć jedną komórkę, nie
zmieniając $\pi_1$; pozostaje prezentacja o $n$ generatorach
i $n-1$ relacjach. Dla diagramu węzła bez skrzyżowań
grupa jest bezpośrednio równa $\mathbb Z$.
\end{proof}

\begin{uwaga}[Rola jednej zbędnej relacji Wirtingera]
Redundancja jednej relacji dotyczy spójnego diagramu węzła
i standardowej prezentacji meridionalnej. Nie oznacza,
że wszystkie pozostałe relacje są algebraicznie niezależne
w dowolnej prezentacji grupy: po wykonaniu transformacji
Tietzego można uzyskać inną liczbę generatorów i relacji.
Przy obliczaniu wielomianu Alexandera należy konsekwentnie
wybierać tę samą abelianizację wszystkich meridianów.
\end{uwaga}

\begin{przyklad}[Trójlistnik]
Dla standardowego trójskrzyżowaniowego diagramu trójlistnika
prezentacja Wirtingera upraszcza się do
$G_T\cong\langle a,b\mid aba=bab\rangle$.
Warto rozwinąć algebraicznie przejście do bardziej symetrycznej
prezentacji
$\langle x,y\mid x^2=y^3\rangle$.
Kładziemy $x:=aba$ oraz $y:=ab$. Z relacji $aba=bab$ wynika
\[
x^2=(aba)(bab)=(ab)^3=y^3.
\]
W drugą stronę, zakładając prezentację $\langle x,y\mid x^2=y^3\rangle$,
definiujemy
$a:=y^{-1}x$ oraz $b:=x^{-1}y^2$. Wtedy $ab=y$, zatem
$aba=(ab)a=ya=x$, a także
$bab=(x^{-1}y^2)(y^{-1}x)(x^{-1}y^2)=x^{-1}y^3=x$.
Stąd $aba=bab$. Otrzymaliśmy wzajemnie odwrotne podstawienia,
więc obie prezentacje opisują tę samą grupę.

Dodatkowo element $z:=x^2=y^3$ jest centralny, gdyż komutuje
zarówno z $x$, jak i z $y$. Po przejściu do ilorazu przez normalną
podgrupę generowaną przez $z$ dostajemy
$\langle x,y\mid x^2=y^3=1\rangle\cong C_2*C_3$.
Ta grupa jest nieprzemienna, a więc i sama grupa trójlistnika
jest nieprzemienna. W szczególności nie jest izomorficzna z
$G_U\cong\mathbb Z$, co ponownie dowodzi nietrywialności
trójlistnika.
\end{przyklad}

\begin{stwierdzenie}
Dla każdego węzła $K\subset S^3$ abelianizacja
$G_K^{\mathrm{ab}}:=G_K/[G_K,G_K]$ jest izomorficzna z $\mathbb Z$.
W szczególności wszystkie meridiany w grupie węzła mają
jednakowy obraz w abelianizacji (po ustaleniu orientacji).
\end{stwierdzenie}

\begin{proof}
Twierdzenie Hurewicza w stopniu pierwszym daje
$G_K^{\mathrm{ab}}\cong H_1(S^3\setminus K;\mathbb Z)$.
Poprzedni rezultat o homologii uzupełnienia splotu zastosowany
do jednej składowej daje $H_1\cong\mathbb Z$, generowaną
przez meridian.
\end{proof}

\begin{uwaga}[Dlaczego sama grupa nie wystarcza]
Grupa $G_K$ jest niezwykle silnym niezmiennikiem, jednak do
odtworzenia pełnej struktury uzupełnienia naturalnie dodaje się
dane \emph{peryferyjne}: wyróżnioną klasę sprzężoności meridianu
i długości (longitude). Przykładowo różne węzły torusowe mogą
mieć grupy abstrakcyjnie powiązane, a badanie włożenia podgrupy
peryferyjnej oddziela informacje gubione przy samej abelianizacji.
\end{uwaga}

\section{Kolorowania Foxa i wyznacznik węzła}
\label{sec:kt-fox}

Kolorowania Foxa stanowią pierwszy przykład niezmiennika, który
można obliczyć metodami elementarnej algebry liniowej.

\begin{definicja}[Kolorowanie $n$-kolorowe]
Niech $n\geq2$ będzie liczbą całkowitą, a $D$ diagramem węzła
$K\subset S^3$. \emph{$n$-kolorowaniem Foxa} nazywamy przypisanie
każdemu łukowi Wirtingera elementu $\mathbb Z/n\mathbb Z$ takiego,
że przy każdym skrzyżowaniu zachodzi
$2b\equiv a+c\pmod n$, gdzie $b$ jest kolorem pasma górnego,
a $a,c$ kolorami dwóch fragmentów pasma dolnego.
Kolorowanie jest \emph{trywialne}, gdy wszystkie łuki mają ten sam kolor.
\end{definicja}

Każdy diagram dopuszcza dokładnie $n$ kolorowań stałych.
Kolorowań nietrywialnych może nie być. Relacje są liniowe,
zatem ich zbiór jest jądrem macierzy nad $\mathbb Z/n\mathbb Z$.

\begin{stwierdzenie}
Dla każdego całkowitego $n\geq2$ liczba $n$-kolorowań Foxa diagramu
węzła nie zależy od diagramu i jest niezmiennikiem węzła.
\end{stwierdzenie}


\begin{proof}
Rozważmy lokalnie trzy ruchy Reidemeistera. W ruchu $\mathrm{R I}$
dodany łuk $b$ ma kolor $a$, gdyż $2a=a+b$ implikuje $b=a$.
W ruchu $\mathrm{R II}$ kolor nowego łuku jest jednoznacznie
wyznaczony przez dwa stare kolory; druga relacja odzyskuje
kolor wyjściowy.
W ruchu $\mathrm{R III}$ dwie kolejności propagowania kolorów
prowadzą do tej samej wartości: definiując $R_b(a)=2b-a$,
obliczamy $R_c(R_b(a))=a+2(c-b)$ oraz
$R_{R_c(b)}(R_c(a))=a+2(c-b)$. Są to dokładnie dwie
możliwe kolejności przejścia przez trzy pasma. Każde $R_b$
jest inwolucją, ponieważ $R_b(R_b(a))=a$; nigdzie nie
trzeba dzielić przez dwa. W ten sposób ruchy indukują
wzajemnie odwrotne bijekcje zbiorów kolorowań.
\end{proof}

\begin{przyklad}[Trójlistnik jest nietrywialny]
Niech $a,b,c\in\mathbb Z/3\mathbb Z$ będą kolorami trzech łuków
standardowego diagramu trójlistnika. Relacje przy skrzyżowaniach
mają postać $2a=b+c$, $2b=a+c$, $2c=a+b$.
W $\mathbb Z/3\mathbb Z$ każda z nich jest równoważna warunkowi
$a+b+c=0$. Istnieje zatem $3^2=9$ kolorowań, z których
trzy są stałe. Węzeł trywialny ma tylko trzy kolorowania.
Zatem trójlistnik nie jest równoważny węzłowi trywialnemu.
\end{przyklad}

\begin{definicja}[Wyznacznik węzła]
Dla węzła $K\subset S^3$ wybieramy diagram $D$ o $n$ łukach
i $n$ skrzyżowaniach. Niech $A_D$ będzie całkowitą macierzą
relacji $2b-a-c=0$, gdzie kolumny odpowiadają łukom, a wiersze
skrzyżowaniom. Po skreśleniu dowolnego jednego wiersza i jednej
kolumny otrzymujemy macierz $(n-1)\times(n-1)$, której wyznacznik
co do wartości bezwzględnej oznaczamy $\det(K)$.
Dla węzła trywialnego przyjmujemy $\det(U)=1$.
\end{definicja}

Niezależność od skreślonego wiersza i kolumny oraz od diagramu
można wykazać, wykorzystując związki macierzy Foxa z wielomianem
Alexandera. Otrzymamy poniżej wzór $\det(K)=|\Delta_K(-1)|$.

\begin{przyklad}
Dla trójlistnika macierz relacji może mieć postać
\[
A_D=\begin{pmatrix}2&-1&-1\\-1&2&-1\\-1&-1&2\end{pmatrix}.
\]
Wyznacznik minora $\begin{psmallmatrix}2&-1\\-1&2\end{psmallmatrix}$
wynosi $3$, zatem $\det(T)=3$. Dla węzła ósemkowego
analogiczny rachunek daje $\det(4_1)=5$.
\end{przyklad}

\begin{stwierdzenie}[Kolorowania a wyznacznik]
Niech $K\subset S^3$ będzie węzłem, $n\ge2$, a
$\det(K)$ jego wyznacznikiem. Jeśli
$\gcd(n,\det(K))=1$, to każde $n$-kolorowanie Foxa
jest stałe. Gdy liczba pierwsza $p$ dzieli $\det(K)$,
istnieje nietrywialne $p$-kolorowanie.
\end{stwierdzenie}

\begin{proof}
Niech $A_D$ będzie macierzą relacji kolorowania, a
$B$ otrzymamy, usuwając jeden wiersz i jedną kolumnę.
Usunięta relacja jest zależna od pozostałych, natomiast
usuniętą współrzędną można ustawić na zero, odejmując
od wszystkich kolorów wspólną stałą. Zatem kolorowania
modulo kolorowania stałe są identyfikowane z
$\ker(B\bmod n)$. Macierz $B$ jest kwadratowa i
$|\det B|=\det(K)$. Jeśli $\gcd(n,\det B)=1$,
macierz $B$ jest odwracalna nad $\mathbb Z/n\mathbb Z$,
więc ma zerowe jądro. Jeśli $p$ dzieli wyznacznik,
$B\bmod p$ jest macierzą osobliwą nad ciałem
$\mathbb F_p$ i posiada niezerowy wektor jądra.
Jest on nietrywialnym kolorowaniem po odtworzeniu
odjętej stałej.
\end{proof}

\section{Powierzchnie Seiferta i rodzaj węzła}
\label{sec:kt-seifert}

\begin{definicja}[Powierzchnia Seiferta]
Niech $K\subset S^3$ będzie zorientowanym węzłem.
\emph{Powierzchnią Seiferta} węzła $K$ nazywamy zwartą,
spójną, zorientowaną, gładko zanurzoną powierzchnię
$F\subset S^3$ taką, że $\partial F=K$ oraz orientacja
brzegu indukowana z $F$ zgadza się z orientacją $K$.
\emph{Rodzajem węzła} jest $g(K):=\min\{g(F):F\text{ jest
powierzchnią Seiferta }K\}$.
\end{definicja}

W pierwszej chwili nie jest oczywiste, że powierzchnia Seiferta
zawsze istnieje. Następne twierdzenie dostarcza konstrukcji.

\begin{twierdzenie}[Seifert; algorytm wygładzania]
Każdy zorientowany węzeł $K\subset S^3$ ogranicza powierzchnię
Seiferta. Jeżeli $D$ jest spójnym diagramem z $c(D)$ skrzyżowaniami,
a po zorientowanym wygładzeniu wszystkich skrzyżowań otrzymujemy
$s(D)$ okręgów Seiferta, to algorytm Seiferta konstruuje spójną
powierzchnię $F_D$ o brzegu $K$ i
\[
\chi(F_D)=s(D)-c(D),\qquad
 g(F_D)=\frac{1-s(D)+c(D)}{2}.
\]
W szczególności $g(K)\leq(1-s(D)+c(D))/2$.
\end{twierdzenie}

\begin{proof}
Przy każdym skrzyżowaniu wygładzamy oba przecinające się łuki
w jedyny sposób zgodny z ich orientacjami. Powstaje skończona
rodzina rozłącznych, zorientowanych okręgów na płaszczyźnie.
Każdy okrąg wypełniamy dyskiem, umieszczając dyski na różnych
wysokościach wtedy, gdy są zagnieżdżone. W otoczeniu każdego
oryginalnego skrzyżowania przytwierdzamy do dwóch odpowiednich
dysków skręcone pasmo, które odtwarza lokalny przebieg nad i pod.
Powstaje gładko zanurzona, zorientowana powierzchnia $F_D$.
Zgodność orientacji na jej fragmentach wynika z zastosowanego
zorientowanego wygładzenia. Jej brzeg jest izotopijny z $K$.
Ponieważ diagram węzła jest spójny jako graf, układ dysków
i pasm również jest spójny. Zaczynając od $s(D)$ dysków
i dołączając $c(D)$ jednowymiarowych uchwytów, otrzymujemy
$\chi(F_D)=s(D)-c(D)$. Dla spójnej orientowalnej powierzchni
z jednym brzegiem $\chi(F_D)=1-2g(F_D)$, co daje podany wzór.
\end{proof}

\begin{center}
\begin{tikzpicture}[scale=.95,>=Stealth,line cap=round,line join=round]
  \node[font=\small\bfseries,text=ktInk] at (-3,1.27) {skrzyżowanie};
  \node[font=\small\bfseries,text=ktInk] at (2.4,1.27) {wygładzenie orientowane};
  \begin{scope}[shift={(-3,0)}]
    \draw[ktBlue,line width=2.9pt,-{Stealth[length=2.7mm]}] (-.9,-.9)--(.9,.9);
    \draw[ktCoral,line width=2.9pt] (-.9,.9)--(-.14,.14);
    \draw[ktCoral,line width=2.9pt,-{Stealth[length=2.7mm]}] (.14,-.14)--(.9,-.9);
  \end{scope}
  \draw[->,very thick,ktTeal] (-1.3,0)--(.25,0);
  \node[text=ktTeal,font=\footnotesize] at (-.53,.3) {wygładź};
  \begin{scope}[shift={(2.4,0)}]
    \draw[ktBlue,line width=2.9pt,-{Stealth[length=2.7mm]}]
       (-.85,-.9) .. controls (-.85,.05) and (-.3,.22) .. (-.3,.9);
    \draw[ktCoral,line width=2.9pt,-{Stealth[length=2.7mm]}]
       (.85,.9) .. controls (.85,-.05) and (.3,-.22) .. (.3,-.9);
  \end{scope}
  \node[font=\footnotesize,text=ktInk] at (0,-1.5)
    {Po wygładzeniu nie ma skrzyżowania; pasma służą później do odbudowania powierzchni.};
\end{tikzpicture}
\end{center}


\begin{center}
\begin{tikzpicture}[scale=.9,line join=round,line cap=round]
  \fill[ktBlue!16,draw=ktBlue,line width=1.4pt] (-2.15,0) ellipse (1.00 and .72);
  \fill[ktCoral!16,draw=ktCoral,line width=1.4pt] (2.15,0) ellipse (1.00 and .72);
  \fill[ktPurple!24,draw=ktPurple,line width=1.25pt]
    (-1.18,.30)
      .. controls (-.72,.34) and (-.28,.48) .. (-.04,.28)
      .. controls (.23,.05) and (.23,-.05) .. (-.04,-.28)
      .. controls (-.28,-.48) and (-.72,-.34) .. (-1.18,-.30)
      .. controls (-.25,-.44) and (.25,.44) .. (1.18,.30)
      .. controls (.72,.34) and (.28,.48) .. (.04,.28)
      .. controls (-.23,.05) and (-.23,-.05) .. (.04,-.28)
      .. controls (.28,-.48) and (.72,-.34) .. (1.18,-.30)
      .. controls (.25,-.44) and (-.25,.44) .. cycle;
  \draw[ktPurple,line width=1.85pt]
    (-1.18,.30) .. controls (-.25,.44) and (.25,-.44) .. (1.18,-.30);
  \draw[ktPurple,line width=1.85pt]
    (-1.18,-.30) .. controls (-.25,-.44) and (.25,.44) .. (1.18,.30);
  \node[text=ktBlue,font=\small] at (-2.15,-1.08) {dysk};
  \node[text=ktCoral,font=\small] at (2.15,-1.08) {dysk};
  \draw[ktPurple,->,thick] (0,-1.08)--(.06,-.42);
  \node[text=ktPurple,font=\small] at (0,-1.32) {pasmo};
  \node[text=ktInk,font=\footnotesize] at (0,-1.80)
    {Poglądowy model dysków Seiferta połączonych pasmem o gładnym styku na brzegach.};
\end{tikzpicture}
\end{center}

\begin{przyklad}
Standardowy diagram trójlistnika ma $c(D)=3$ i $s(D)=2$,
więc algorytm daje $g(F_D)=1$. Standardowy diagram węzła
ósemkowego ma $c(D)=4$, $s(D)=3$, a więc również $g(F_D)=1$.
Dopiero obliczenie niezmienników algebraicznych wykaże, że
żaden z tych węzłów nie ogranicza dysku; ich rodzaj jest
rzeczywiście równy jeden.
\end{przyklad}

\begin{stwierdzenie}
Dla węzła $K\subset S^3$ następujące warunki są równoważne:
(a) $K$ jest trywialny; (b) $g(K)=0$; (c) $K$ ogranicza
gładko zanurzony dysk w $S^3$.
\end{stwierdzenie}

\begin{proof}
Równoważność (b) i (c) wynika z klasyfikacji zwartych,
spójnych, orientowalnych powierzchni z jednym brzegiem:
powierzchnia rodzaju zero jest dyskiem. Jeżeli $K$ ogranicza
dysk, jego otoczenie regularne w $S^3$ jest kulą, a brzeg
dysku można wyprostować izotopią otoczenia do okręgu standardowego.
W drugą stronę standardowy okrąg ogranicza dysk, którego obraz
przy izotopii ogranicza $K$.
\end{proof}

\begin{definicja}[Gładki rodzaj czterowymiarowy i węzeł typu slice]
\label{def:kt-four-ball-genus}
Niech $K\subset S^3=\partial B^4$ będzie węzłem.
Przez $g_4^{\mathrm{sm}}(K)$ oznaczamy najmniejszy rodzaj
zwartej, spójnej, zorientowanej, gładko i właściwie
zanurzonej powierzchni $F\subset B^4$ z jednym brzegiem
$\partial F=K$. Przy brzegu żądamy standardowego
kołnierza produktu. Węzeł nazywamy
\emph{gładko typu slice}, jeżeli
$g_4^{\mathrm{sm}}(K)=0$, czyli gdy ogranicza
gładko i właściwie zanurzony dysk w $B^4$.
Przesunięcie powierzchni Seiferta do kołnierza
$S^3\times[0,\varepsilon)\subset B^4$
pokazuje, że $g_4^{\mathrm{sm}}(K)\le g(K)$.
\end{definicja}

\section{Macierz Seiferta i wielomian Alexandera}
\label{sec:kt-alexander}

Niech $F\subset S^3$ będzie powierzchnią Seiferta węzła $K$,
o rodzaju $g$. Grupa $H_1(F;\mathbb Z)$ jest wolna rzędu $2g$.
Niech $[\alpha_1],\dots,[\alpha_{2g}]$ będzie jej bazą,
reprezentowaną przez zamknięte krzywe na $F$. Ustalmy dodatnią
normalną do $F$; przez $\alpha_i^+$ oznaczamy nieznaczne
przesunięcie $\alpha_i$ w tym kierunku, rozłączne z $F$.

\begin{definicja}[Macierz Seiferta]
Macierzą Seiferta powierzchni $F$ w wybranej bazie nazywamy
macierz całkowitą $V_F=(v_{ij})_{i,j=1}^{2g}$ daną wzorem
$v_{ij}:=\operatorname{lk}(\alpha_i^+,\alpha_j)$.
Jej symetryzacją jest $V_F+V_F^T$.
\end{definicja}

\begin{stwierdzenie}[Postać przecięcia]
Niech $F$ będzie zorientowaną powierzchnią Seiferta węzła
$K\subset S^3$, a $V=V_F$ jej macierzą Seiferta. Wtedy
$V-V^T$ reprezentuje, z dokładnością do znaku wynikającego
z konwencji normalnej, algebraiczną postać przecięcia
na $H_1(F;\mathbb Z)$. W szczególności $\det(V-V^T)=1$.
\end{stwierdzenie}

\begin{proof}
Różnica $\operatorname{lk}(\alpha_i^+,\alpha_j)
-\operatorname{lk}(\alpha_j^+,\alpha_i)$ równa się algebraicznej
liczbie przecięć krzywych $\alpha_i$ i $\alpha_j$ na powierzchni
(z ewentualną wspólną zmianą znaku zależnie od orientacji).
Jest to lokalny rachunek: gdy dwie krzywe przecinają się
poprzecznie, odpowiednie dwa przesunięcia tworzą jedno
skrzyżowanie liczone z właściwym znakiem. Dla spójnej
orientowanej powierzchni rodzaju $g$ z jednym brzegiem
postać przecięcia ma w bazie symplektycznej macierz
$\begin{psmallmatrix}0&I_g\\-I_g&0\end{psmallmatrix}$,
której wyznacznik wynosi jeden. Zmiana bazy w $GL_{2g}(\mathbb Z)$
nie zmienia wyznacznika.
\end{proof}

\begin{definicja}[Wielomian Alexandera]
Dla węzła $K\subset S^3$ o powierzchni Seiferta $F$
i macierzy $V$ definiujemy
$\Delta_K(t):=\det(tV-V^T)\in\mathbb Z[t,t^{-1}]$,
z dokładnością do mnożenia przez jednostkę $\pm t^k$.
Dla powierzchni będącej dyskiem przyjmujemy wyznacznik pustej
macierzy równy jeden. Jeśli zachodzi taka potrzeba, wybieramy
symetryczną normalizację spełniającą
$\Delta_K(t)=\Delta_K(t^{-1})$ oraz $\Delta_K(1)=1$.
\end{definicja}

\begin{uwaga}
W definicji należy dopiero wykazać, że wynik nie zależy
od powierzchni Seiferta ani od jej bazy. Wyznacznik $\det(tV-V^T)$
zmienia się pod zmianą bazy przez mnożenie po lewej i prawej
przez odpowiednie macierze unimodularne, więc nie zmienia
się wcale. Niezależność od powierzchni wymaga dodatkowej idei:
równoważności $S$ macierzy Seiferta lub modułu Alexandera.
\end{uwaga}

\begin{twierdzenie}[Opis przez nakrycie cykliczne]
Niech $K\subset S^3$ będzie węzłem, $X_K$ jego uzupełnieniem,
a $\widetilde X_K\to X_K$ nieskończonym nakryciem cyklicznym
wyznaczonym przez epimorfizm
$\pi_1(X_K)\to H_1(X_K;\mathbb Z)\cong\mathbb Z$.
Grupa $H_1(\widetilde X_K;\mathbb Z)$ jest modułem nad
$\Lambda:=\mathbb Z[t,t^{-1}]$, przy czym $t$ działa przez
wybrany generator grupy przekształceń pokładowych. Dla
macierzy Seiferta $V$ moduł ten posiada kwadratową macierz
prezentacji $tV-V^T$. Stąd $\Delta_K(t)$ jest niezmiennikiem
węzła z dokładnością do $\pm t^k$.
\end{twierdzenie}

\begin{proof}[Rozcięcie, nakrycie i ciąg Mayera--Vietorisa]
Niech $F$ będzie powierzchnią Seiferta rodzaju $g$, a
$Y$ domknięciem $X_K$ rozciętego wzdłuż właściwie
zanurzonej części $F\cap X_K$. W brzegu $Y$ znajdują się
dwie kopie tej powierzchni, oznaczone $F_+$ i $F_-$.
Nieskończone nakrycie cykliczne $\widetilde X_K$ otrzymujemy,
klejąc kopie $Y_j$, $j\in\mathbb Z$, tak, by strona
$F_+$ kopii $Y_j$ była utożsamiona ze stroną $F_-$
kopii $Y_{j+1}$. Działanie $t$ przesuwa numer $j$ o jeden.

Zapiszmy ciąg Mayera--Vietorisa nad
$\Lambda=\mathbb Z[t,t^{-1}]$:
\[
H_1(F)\otimes\Lambda\xrightarrow{\;t(i_+)_*-(i_-)_*\;}
H_1(Y)\otimes\Lambda\longrightarrow
H_1(\widetilde X_K)\longrightarrow
H_0(F)\otimes\Lambda\xrightarrow{\;t-1\;}
H_0(Y)\otimes\Lambda.
\]
Powierzchnia $F$ i rozcięta przestrzeń $Y$ są spójne.
Ponieważ $\Lambda$ jest dziedziną całkowitości, mnożenie
przez $t-1$ jest injekcją. Z dokładności otrzymujemy
$H_1(\widetilde X_K)\cong\operatorname{coker}
(t(i_+)_*-(i_-)_*)$.

Aby obliczyć macierz tej mapy, zauważmy, że $Y$ jest
homologicznie równoważne dopełnieniu powierzchni Seiferta.
Dualność Alexandera daje
$H_1(Y;\mathbb Z)\cong H^1(F;\mathbb Z)\cong\mathbb Z^{2g}$.
Wybierzmy bazę $[\alpha_1],\dots,[\alpha_{2g}]$ w $H_1(F)$,
a w $H_1(Y)$ bazę dualną, mierzoną za pomocą liczby splecenia
z tymi krzywymi. Macierze dwóch włożeń są wtedy macierzami
formy Seiferta $V$ i $V^T$ (zależnie od ustalenia numeracji
stron). Stąd macierz homomorfizmu wynosi $tV-V^T$.

Zmiana bazy nie zmienia modułu zdefiniowanego topologicznie.
Ponieważ $\det(V-V^T)=1$, wyznacznik $\det(tV-V^T)$
nie jest wielomianem zerowym. Zerowy ideał Fittinga
tego torsyjnego modułu jest zatem generowany przez
$\det(tV-V^T)$; inne bazy i powierzchnie dają ten sam
ideał. Jednostkami $\Lambda$ są dokładnie $\pm t^k$,
co kończy dowód niezależności wielomianu.
\end{proof}

\begin{stwierdzenie}[Własności wielomianu Alexandera]
Dla każdego węzła $K\subset S^3$ wielomian $\Delta_K$ spełnia:
\begin{enumerate}[label=(\roman*)]
\item $\Delta_K(1)=\pm1$ dla dowolnego reprezentanta, a w normalizacji
symetrycznej $\Delta_K(1)=1$;
\item $\Delta_K(t^{-1})=\pm t^m\Delta_K(t)$ dla pewnego $m\in\mathbb Z$;
\item $\det(K)=|\Delta_K(-1)|$, stąd $\det(K)$ jest dodatnią liczbą nieparzystą;
\item szerokość wykładników $\Delta_K$, czyli
$\deg_{\max}\Delta_K-\deg_{\min}\Delta_K$, nie przekracza $2g(K)$.
\end{enumerate}
\end{stwierdzenie}

\begin{proof}
Z twierdzenia o postaci przecięcia
$\Delta_K(1)=\det(V-V^T)=1$ dla przyjętej reprezentacji.
Przy $2g\times2g$ otrzymujemy
$\det(t^{-1}V-V^T)=t^{-2g}\det(V-tV^T)
=t^{-2g}\det(tV-V^T)$, gdyż wyznacznik jest niezmienniczy
przy transpozycji. Daje to symetrię. Wzór na wyznacznik
węzła wynika z identyfikacji macierzy Foxa przy $t=-1$
z odpowiednią prezentacją modułu Alexandera; jest to standardowa
redukcja relacji Wirtingera. Nieparzystość wynika z kongruencji
$\Delta_K(-1)\equiv\Delta_K(1)\equiv1\pmod2$.
Wreszcie macierz $tV-V^T$ ma rozmiar $2g$, więc jej
wyznacznik jest wielomianem stopnia co najwyżej $2g$
z nieujemnymi potęgami $t$. Minimalizując $g$ po wszystkich
powierzchniach, uzyskujemy ostatnią nierówność.
\end{proof}

\begin{przyklad}[Dwa najprostsze nietrywialne węzły]
Dla jednego z odbić trójlistnika można wybrać macierz
$V_T=\begin{psmallmatrix}1&0\\-1&1\end{psmallmatrix}$. Wtedy
\[
tV_T-V_T^T=
\begin{pmatrix}t-1&1\\-t&t-1\end{pmatrix},
\qquad
\Delta_T(t)=\det(tV_T-V_T^T)=(t-1)^2+t=t^2-t+1.
\]
Po symetrycznej normalizacji otrzymujemy
$\Delta_T^{\mathrm{sym}}(t)=t^{-1}-1+t$.
Ponieważ szerokość wykładników wynosi $2$, dostajemy
$g(T)\ge 1$. Z algorytmu Seiferta wiemy już, że $g(T)\le 1$,
a więc $g(T)=1$. Wyznacznik węzła jest równy
$\det(T)=|\Delta_T(-1)|=3$.

Ta sama macierz pozwala natychmiast odczytać sygnaturę.
Mamy bowiem
\[
V_T+V_T^T=
\begin{pmatrix}2&-1\\-1&2\end{pmatrix},
\]
której wartości własne wynoszą $1$ i $3$. Są dodatnie,
zatem $\sigma(T)=2$. W szczególności trójlistnik nie jest
amfichiralny, a z klasycznej nierówności sygnaturowej,
sformułowanej w Twierdzeniu~\ref{thm:kt-signature-four-genus},
otrzymujemy $g_4^{\mathrm{sm}}(T)\ge 1$.

Dla węzła ósemkowego można wziąć macierz
$V_8=\begin{psmallmatrix}1&0\\-1&-1\end{psmallmatrix}$. Otrzymujemy
\[
tV_8-V_8^T=
\begin{pmatrix}t-1&1\\-t&1-t\end{pmatrix},
\qquad
\Delta_{4_1}(t)=\det(tV_8-V_8^T)=-t^2+3t-1.
\]
Po normalizacji
$\Delta_{4_1}^{\mathrm{sym}}(t)=3-t-t^{-1}$,
stąd ponownie $g(4_1)=1$ oraz
$\det(4_1)=|\Delta_{4_1}(-1)|=5$.
Z kolei
\[
V_8+V_8^T=
\begin{pmatrix}2&-1\\-1&-2\end{pmatrix}
\]
ma wyznacznik ujemny, a więc jest nieokreślona; stąd
$\sigma(4_1)=0$. Oba węzły mają więc ten sam rodzaj,
ale różnią się wielomianem Alexandera, wyznacznikiem i sygnaturą.
\end{przyklad}

\begin{uwaga}[Ograniczenia wielomianu]
Równość $\Delta_K(t)=1$ nie pociąga za sobą trywialności węzła.
Istnieją nietrywialne węzły o trywialnym wielomianie Alexandera;
przykładowo odpowiednie satelity typu Whiteheada. Wielomian
Alexandera nie rozróżnia też węzła od jego odbicia lustrzanego.
Potrzebujemy zatem dalszych niezmienników.
\end{uwaga}

\begin{stwierdzenie}[Zmiana bazy i stabilizacja macierzy Seiferta]
Niech $K\subset S^3$ będzie węzłem i niech $V$ będzie
jego macierzą Seiferta. Przy zastąpieniu bazy w
$H_1(F;\mathbb Z)$ macierz przechodzi w $P^TVP$,
gdzie $P\in GL_{2g}(\mathbb Z)$. Przy elementarnym
powiększeniu powierzchni Seiferta macierz można, po
odpowiedniej zmianie bazy, zapisać jako
\[
V'=\begin{pmatrix}V&u&0\\ v^T&z&0\\0&1&0\end{pmatrix}
\quad\text{lub}\quad
V'=\begin{pmatrix}V&u&0\\v^T&z&1\\0&0&0\end{pmatrix},
\]
gdzie $u,v\in\mathbb Z^{2g}$ są wektorami kolumnowymi,
a $z\in\mathbb Z$. W obu przypadkach
$\det(tV'-{V'}^T)$ różni się od $\det(tV-V^T)$
co najwyżej przez czynnik będący jednostką w
$\mathbb Z[t,t^{-1}]$.
\end{stwierdzenie}

\begin{proof}
Przy zmianie bazy biliniowa forma Seiferta jest mnożona
przez $P^T$ i $P$, zatem
$\det(tP^TVP-P^TV^TP)=(\det P)^2\det(tV-V^T)
=\det(tV-V^T)$. Dla pierwszej postaci ostatnia kolumna macierzy
$tV'-{V'}^T$ ma jeden niezerowy element $-1$
w przedostatnim wierszu, a jej ostatni wiersz
ma jeden niezerowy element $t$ w przedostatniej
kolumnie. Rozwinięcie wyznacznika wzdłuż tych
dwóch pozycji daje dokładnie
$\det(tV'-{V'}^T)=t\det(tV-V^T)$.
Dla drugiej postaci odpowiednie elementy
wynoszą $t$ oraz $-1$ i otrzymujemy tę samą
równość. Wektory $u,v$ i liczba $z$ nie wpływają
na ten rachunek. Czynnik $t$ jest jednostką
pierścienia $\mathbb Z[t,t^{-1}]$.
Zera w ostatniej kolumnie, odpowiednio
w ostatnim wierszu, są niezbędne: dopuszczenie
tam dowolnych wektorów nie jest ogólnie
elementarnym powiększeniem i może zmienić
wielomian Alexandera poza mnożeniem przez jednostkę.
Twierdzenie o $S$-równoważności stwierdza, że wszystkie macierze
Seiferta tego samego węzła można połączyć skończonym
ciągiem takich operacji i ich odwrotności. Ten ostatni
fakt jest niezależnym twierdzeniem topologicznym; powyższy
rachunek pokazuje, dlaczego wystarcza on do wykazania
niezmienniczości $\Delta_K$.
\end{proof}

\section{Rachunek Foxa i algorytm wyznaczania wielomianu Alexandera}
\label{sec:kt-fox-calculus}

Macierz Seiferta jest geometrycznym sposobem obliczania
wielomianu Alexandera. Jeżeli znamy jednak prezentację
Wirtingera, nie musimy konstruować powierzchni: wystarcza
specjalny rachunek różniczkowy w pierścieniu grupowym.
Metoda ta działa bezpośrednio na słowach w generatorach.

\begin{definicja}[Pochodne Foxa]
Niech $F_n=\langle x_1,\dots,x_n\rangle$ będzie wolną grupą,
a $\mathbb Z[F_n]$ jej całkowitym pierścieniem grupowym.
Dla $i\in\{1,\dots,n\}$ \emph{pochodną Foxa względem $x_i$}
nazywamy jedyne $\mathbb Z$-liniowe odwzorowanie
$\partial_i\colon\mathbb Z[F_n]\to\mathbb Z[F_n]$
spełniające $\partial_i(x_j)=\delta_{ij}$ i
$\partial_i(uv)=\partial_i(u)+u\partial_i(v)$
dla wszystkich $u,v\in F_n$. W szczególności
$\partial_i(1)=0$ oraz
$\partial_i(x_j^{-1})=-x_j^{-1}\delta_{ij}$.
\end{definicja}

Różnicą w stosunku do zwykłej pochodnej jest reguła
Leibniza: drugi składnik jest mnożony przez $u$, a nie
przez pochodną $u$ razy $v$. Operacja ta jest naturalną
pochodną w wolnym rachunku różniczkowym, dopasowaną
do algebry grupy.

\begin{lemat}[Fundamentalna tożsamość Foxa]
Dla każdego słowa $w\in F_n$ zachodzi równość
$w-1=\sum_{i=1}^n\partial_i(w)(x_i-1)$
w pierścieniu $\mathbb Z[F_n]$.
\end{lemat}

\begin{proof}
Dla generatora $w=x_j$ prawa strona wynosi $x_j-1$.
Dla odwrotności $w=x_j^{-1}$ mamy
$-x_j^{-1}(x_j-1)=x_j^{-1}-1$. Jeśli tożsamość
zachodzi dla $u$ i $v$, to dla $uv$ prawa strona wynosi
$\sum_i(\partial_i(u)+u\partial_i(v))(x_i-1)
=(u-1)+u(v-1)=uv-1$.
Indukcja względem długości słowa kończy dowód.
\end{proof}

\begin{definicja}[Macierz Alexandera z prezentacji grupy]
Niech $G=\langle x_1,\dots,x_n\mid r_1,\dots,r_m\rangle$
będzie skończoną prezentacją grupy, a
$\phi\colon G\to\mathbb Z=\langle t\rangle$ epimorfizmem.
Przez ten sam symbol oznaczamy indukowany homomorfizm
pierścieni $\mathbb Z[F_n]\to\mathbb Z[t,t^{-1}]$.
\emph{Macierzą Foxa--Alexandera} tej prezentacji i
homomorfizmu $\phi$ nazywamy macierz
$A=(\phi(\partial_jr_i))_{i=1,\dots,m;\ j=1,\dots,n}$
nad pierścieniem $\Lambda=\mathbb Z[t,t^{-1}]$.
\end{definicja}

\begin{twierdzenie}[Algorytm dla diagramu Wirtingera]
Niech $D$ będzie diagramem zorientowanego węzła
$K\subset S^3$ o $n\geq1$ skrzyżowaniach.
Z prezentacji Wirtingera wybierzmy $n$ generatorów
$x_1,\dots,x_n$ i $n-1$ niezależnych relacji
$r_1,\dots,r_{n-1}$. Niech
$\phi(x_j)=t$ dla wszystkich $j$. Zbudujmy macierz
$A=(\phi(\partial_jr_i))$ rozmiaru $(n-1)\times n$.
Po skreśleniu dowolnej jednej kolumny otrzymujemy
macierz kwadratową $A^{(j)}$, dla której
$\Delta_K(t)\doteq\det A^{(j)}$.
Jeżeli wybrana prezentacja ma dodatkowe generatory
lub relacje, stosujemy ogólną teorię ideałów
wyznacznikowych macierzy Alexandera.
\end{twierdzenie}

\begin{proof}
Niech $\Lambda=\mathbb Z[t,t^{-1}]$.
Standardowy kompleks dwuwymiarowy prezentacji Wirtingera
ma po podniesieniu do nakrycia cyklicznego łańcuchy
$C_2=\Lambda^{n-1}$, $C_1=\Lambda^n$ i $C_0=\Lambda$.
Operator $\partial_1$ jest dany wzorem
$\partial_1(u_1,\dots,u_n)=(t-1)\sum_{j=1}^n u_j$,
gdyż wszystkie meridiany przechodzą na $t$.
Ponieważ $t-1$ nie jest dzielnikiem zera w $\Lambda$,
$\ker\partial_1$ składa się z wektorów o sumie współrzędnych
równej zero. Usunięcie dowolnej współrzędnej daje więc
izomorfizm $\Lambda$-modułów
$\ker\partial_1\cong\Lambda^{n-1}$; brakującą współrzędną
odzyskuje się jako minus sumę pozostałych.

Macierz $\partial_2$, w konwencji kolumnowej, jest
transpozycją macierzy pochodnych Foxa. Istotnie, śledząc
brzeg podniesionej komórki odpowiadającej słowu
$r_i$, każdą literę $x_j$ odczytujemy z wagą równą obrazowi
prefiksu tego słowa w $\Lambda$. Jest to dokładnie reguła
Leibniza pochodnej Foxa; litera odwrotna wnosi wagę ujemną.
Fundamentalna tożsamość Foxa implikuje
$\partial_1\partial_2=0$. Dlatego po utożsamieniu
$\ker\partial_1$ z $\Lambda^{n-1}$ macierz
$\partial_2$ jest transpozycją $A^{(j)}$.

W pierwszej homologii mamy
$H_1(\widetilde X_K;\mathbb Z)=\ker\partial_1/
\operatorname{im}\partial_2$. Standardowy kompleks
Wirtingera odtwarza pierwszy typ homotopii dopełnienia,
więc oblicza tę samą pierwszą homologię co nakrycie.
Zatem moduł Alexandera ma prezentację $A^{(j)}$
(z dokładnością do transpozycji), a jego wielomianem
rzędu zerowego jest $\det A^{(j)}$. W połączeniu z
poprzednim twierdzeniem daje to
$\det A^{(j)}\doteq\Delta_K(t)$.
\end{proof}

\begin{przyklad}[Pełne obliczenie dla trójlistnika]
Weźmy prezentację
$G_T=\langle a,b\mid aba=bab\rangle$.
Relację zapisujemy jako słowo
$r=aba\,b^{-1}a^{-1}b^{-1}$.
Reguła Foxa daje
$\partial_a r=1+ab-aba\,b^{-1}a^{-1}$,
a po abelianizacji $\phi(a)=\phi(b)=t$ otrzymujemy
$\phi(\partial_a r)=1+t^2-t=t^2-t+1$.
Podobnie
$\partial_b r=a-aba\,b^{-1}-aba\,b^{-1}a^{-1}b^{-1}$,
więc $\phi(\partial_b r)=t-t^2-1$.
Macierz Alexandera jest zatem
$A=\begin{pmatrix}t^2-t+1&-(t^2-t+1)\end{pmatrix}$.
Usuwając dowolną kolumnę, otrzymujemy
$\Delta_T(t)\doteq t^2-t+1$, zgodnie z rachunkiem
z macierzy Seiferta. Widzimy, że metoda wymaga tylko
operacji na słowach grupowych i wielomianach Laurenta.
\end{przyklad}

\begin{uwaga}[Procedura obliczeniowa]
Dla praktycznych obliczeń można kolejno:
(1) zakodować orientowany diagram i jego łuki Wirtingera;
(2) zapisać relacje jako słowa w wolnej grupie;
(3) policzyć pochodne Foxa z reguł Leibniza;
(4) zastąpić każdy meridian zmienną $t$;
(5) skreślić kolumnę macierzy i policzyć wyznacznik;
(6) usunąć czynnik $\pm t^k$ przez symetryczną normalizację.
Kroki (3)--(5) są czysto algebraiczne i mogą być
zaimplementowane w systemie algebry komputerowej.
\end{uwaga}

\section{Wielomian Conwaya i niezmiennik Arfa}
\label{sec:kt-conway-arf}

\begin{definicja}[Wielomian Conwaya]
Dla każdego zorientowanego splotu $L\subset S^3$
przez $\nabla_L(z)\in\mathbb Z[z]$ oznaczamy
wielomian Conwaya, jednoznacznie scharakteryzowany
przez warunki $\nabla_U(z)=1$ oraz
$\nabla_{L_+}(z)-\nabla_{L_-}(z)=z\nabla_{L_0}(z)$
dla każdej zorientowanej trójki skein. Dla splotu
rozdzielnego o co najmniej dwóch składowych
$\nabla_L(z)=0$.
\end{definicja}

\begin{twierdzenie}[Conway--Alexander]
Dla dowolnego zorientowanego węzła $K\subset S^3$
wielomian Conwaya spełnia
$\Delta_K^{\mathrm{sym}}(t)=
\nabla_K(t^{1/2}-t^{-1/2})$,
gdzie $\Delta_K^{\mathrm{sym}}$ jest symetrycznie
znormalizowanym wielomianem Alexandera.
Ponadto $\nabla_K(z)$ zawiera wyłącznie parzyste
potęgi $z$, a jego wyraz stały wynosi jeden.
\end{twierdzenie}

\begin{proof}[Zarys indukcji skein]
Dla splotów wprowadzamy symetrycznie unormowany wielomian
Alexandera z uwzględnieniem liczby składowych; jego wartości
mogą należeć do $\mathbb Z[t^{1/2},t^{-1/2}]$. Relacja
Alexandera--Conwaya ma postać
$\Delta_{L_+}(t)-\Delta_{L_-}(t)=
(t^{1/2}-t^{-1/2})\Delta_{L_0}(t)$.
Można ją uzyskać bezpośrednio z dwóch macierzy Seiferta
różniących się lokalnym uchwytem w otoczeniu skrzyżowania
lub z odpowiednich macierzy prezentacji modułu Alexandera.
Po podstawieniu $z=t^{1/2}-t^{-1/2}$ przyjmuje ona dokładnie
postać relacji Conwaya.

Wyjaśnijmy jednoznaczność rekurencji. Dla diagramu splotu
wybieramy punkty bazowe na składowych i kolejność tych
składowych. Zmieniając skrzyżowania, w których po raz
pierwszy przechodzimy dołem podczas śledzenia diagramu,
uzyskujemy diagram splotu rozdzielnego. Każdą taką zmianę
wyrażamy relacją skein, przy czym wygładzenie zmniejsza
liczbę skrzyżowań. Indukcja po liczbie skrzyżowań, wraz
z wartością $1$ na węźle trywialnym i $0$ na splotach
rozdzielnych o co najmniej dwóch składowych, wyznacza
wielomian jednoznacznie. Dla węzła relacja symetrii
wielomianu Alexandera odpowiada niezmienniczości po
zamianie $z$ na $-z$, skąd znikają potęgi nieparzyste.
\end{proof}

\begin{definicja}[Forma kwadratowa i niezmiennik Arfa]
Niech $K\subset S^3$ będzie węzłem i niech
$F$ będzie jego powierzchnią Seiferta rodzaju $g$.
Na $H_1(F;\mathbb Z/2\mathbb Z)$ definiujemy
formę kwadratową $q_F([\gamma]):=
\operatorname{lk}(\gamma,\gamma^+)\pmod2$.
Jeżeli $(a_1,b_1,\dots,a_g,b_g)$ jest bazą symplektyczną
przestrzeni $H_1(F;\mathbb Z/2\mathbb Z)$ względem
modulo-dwójkowej formy przecięcia, kładziemy
$\operatorname{Arf}(K):=\sum_{i=1}^{g}q_F(a_i)q_F(b_i)
\pmod2\in\mathbb Z/2\mathbb Z$.
Niezależność od wyborów wykazujemy w następnym stwierdzeniu.
\end{definicja}

\begin{lemat}[Suma Gaussa dla formy kwadratowej nad $\mathbb F_2$]
\label{lem:tg-arf-gauss}
Niech $V$ będzie przestrzenią liniową wymiaru $2g$ nad
$\mathbb F_2$, a $\omega\colon V\times V\to\mathbb F_2$
niezdegenerowaną formą naprzemienną. Niech funkcja
$q\colon V\to\mathbb F_2$ spełnia
$q(x+y)=q(x)+q(y)+\omega(x,y)$ dla wszystkich $x,y\in V$.
Wówczas dla każdej bazy symplektycznej
$(a_1,b_1,\ldots,a_g,b_g)$ zachodzi
\[
\sum_{x\in V}(-1)^{q(x)}
=2^g(-1)^{\sum_{i=1}^g q(a_i)q(b_i)}.
\]
W szczególności suma $\sum_iq(a_i)q(b_i)\in\mathbb F_2$
nie zależy od wyboru bazy symplektycznej.
\end{lemat}

\begin{proof}
Każdy $x\in V$ ma jednoznaczny zapis
$x=\sum_i(\alpha_i a_i+\beta_i b_i)$, gdzie
$\alpha_i,\beta_i\in\mathbb F_2$. Ponieważ różne pary
$(a_i,b_i)$ są względem $\omega$ ortogonalne, a
$\omega(a_i,b_i)=1$, to z tożsamości kwadratowej wynika
\[
q(x)=\sum_{i=1}^g
\bigl(\alpha_i q(a_i)+\beta_i q(b_i)+\alpha_i\beta_i\bigr)
\pmod 2.
\]
Suma Gaussa jest zatem iloczynem $g$ sum dwuwymiarowych.
Dla $A,B\in\mathbb F_2$ mamy
\[
\sum_{\alpha,\beta\in\mathbb F_2}
(-1)^{\alpha A+\beta B+\alpha\beta}
=1+(-1)^A+(-1)^B-(-1)^{A+B}
=2(-1)^{AB}.
\]
Mnożąc te równości, dostajemy wzór. Lewa strona zależy
wyłącznie od $q$, a zatem wyznacza niezależnie od bazy
wykładnik po prawej stronie.
\end{proof}

\begin{stwierdzenie}[Niezmienniczość formy Arfa]
Dla każdego węzła $K\subset S^3$ i powierzchni Seiferta
$F$ funkcja $q_F$ spełnia
$q_F(x+y)=q_F(x)+q_F(y)+x\cdot y$ w $\mathbb F_2$,
gdzie $x\cdot y$ oznacza formę przecięcia na $H_1(F;\mathbb F_2)$.
Liczba $\sum_iq_F(a_i)q_F(b_i)$ jest niezależna od bazy
symplektycznej i powierzchni $F$.
\end{stwierdzenie}

\begin{proof}
Dla dwóch poprzecznych krzywych $\alpha,\beta\subset F$
reprezentant klasy $[\alpha]+[\beta]$ otrzymujemy przez
lokalne wygładzenie punktów przecięcia. Przy obliczaniu
samosplecenia pojawiają się składniki diagonalne
$q_F([\alpha])$, $q_F([\beta])$ oraz suma dwóch składników
mieszanych. Modulo dwa te ostatnie są równe algebraicznej
liczbie przecięć krzywych, stąd wzór kwadratowy.

Niezależność sumy $\sum_iq_F(a_i)q_F(b_i)$ od bazy
symplektycznej wynika dokładnie z lematu
\ref{lem:tg-arf-gauss}: jego lewa strona jest określona
bez użycia jakiejkolwiek bazy.
Przy zmianie powierzchni Seiferta macierze są związane
$S$-równoważnością. Elementarne powiększenie dodaje parę
hiperboliczną, dla której składnik formy Arfa wynosi zero.
Zmiany bazy i takie powiększenia nie zmieniają więc wyniku.
\end{proof}

\begin{stwierdzenie}
Niech $K\subset S^3$ będzie węzłem, a
$\nabla_K(z)=1+a_2z^2+a_4z^4+\cdots$.
Wtedy $\operatorname{Arf}(K)\equiv a_2\pmod2$.
Ponadto jeśli $K$ jest gładko typu slice
w sensie Definicji~\ref{def:kt-four-ball-genus}, to
$\operatorname{Arf}(K)=0$.
\end{stwierdzenie}

\begin{uwaga}
Zerowanie niezmiennika Arfa jest warunkiem koniecznym
slicowości, ale nie wystarczającym. Niezmiennik ten
jest znacznie słabszy niż $\tau$ czy $s$, lecz
ma ważne znaczenie w teorii form kwadratowych
nad ciałem dwuelementowym i w czterowymiarowej
teorii chirurgii.
\end{uwaga}

\begin{przyklad}
Dla trójlistnika
$\Delta_T^{\mathrm{sym}}(t)=t+t^{-1}-1$.
Ponieważ
$(t^{1/2}-t^{-1/2})^2=t+t^{-1}-2$,
mamy $\nabla_T(z)=1+z^2$ i
$\operatorname{Arf}(T)=1$.
Dla węzła ósemkowego
$\Delta_{4_1}^{\mathrm{sym}}(t)=3-t-t^{-1}$,
więc $\nabla_{4_1}(z)=1-z^2$;
także $\operatorname{Arf}(4_1)=1$.
\end{przyklad}


\section{Sygnatura węzła i ograniczenia rodzaju}
\label{sec:kt-sygnatura}

\begin{definicja}[Sygnatura klasyczna]
Dla węzła $K\subset S^3$ i macierzy Seiferta $V$ definiujemy
$\sigma(K):=\operatorname{sign}(V+V^T)$,
gdzie $\operatorname{sign}$ oznacza liczbę dodatnich minus liczbę
ujemnych wartości własnych rzeczywistej macierzy symetrycznej,
liczonych z krotnościami. Wartość nie zależy od wyboru
powierzchni i bazy.
\end{definicja}

Niezależność od powierzchni wynika z $S$-równoważności macierzy
Seiferta: zmiana bazy zachowuje sygnaturę przez prawo
bezwładności Sylvestera, a elementarne powiększenia odpowiadają
dodaniu składnika hiperbolicznego o sygnaturze zero.

\begin{stwierdzenie}
Niech $K\subset S^3$ będzie węzłem. Wówczas
$|\sigma(K)|\leq2g(K)$.
Odbicie lustrzane zmienia znak sygnatury:
$\sigma(\overline K)=-\sigma(K)$.
\end{stwierdzenie}

\begin{proof}
Macierz $V+V^T$ ma rozmiar $2g(F)$, a sygnatura symetrycznej
macierzy wymiaru $2g(F)$ ma wartość bezwzględną nie większą
od $2g(F)$. Minimalizując po powierzchniach, dostajemy pierwszą
nierówność. Przy odbiciu lustrzanym macierz Seiferta zmienia się,
z dokładnością do odpowiedniej zmiany bazy, na $-V^T$;
symetryzacja zmienia znak, co daje drugą równość.
\end{proof}

\begin{przyklad}
Dla $V_T=\begin{psmallmatrix}1&0\\-1&1\end{psmallmatrix}$
macierz $V_T+V_T^T=\begin{psmallmatrix}2&-1\\-1&2\end{psmallmatrix}$
ma dwie dodatnie wartości własne, więc $\sigma(T)=2$
w tej konwencji wyboru odbicia trójlistnika. Jego odbicie
ma sygnaturę $-2$. W szczególności trójlistnik nie jest
amfichiralny. Dla $V_8$ macierz symetryczna ma jedną wartość
dodatnią i jedną ujemną, a więc $\sigma(4_1)=0$.
\end{przyklad}

\begin{twierdzenie}[Nierówność sygnaturowa dla rodzaju czterowymiarowego;
wynik przyjmowany]
\label{thm:kt-signature-four-genus}
Dla każdego węzła $K\subset S^3$ zachodzi
\[
|\sigma(K)|\le 2g_4^{\mathrm{sm}}(K).
\]
W szczególności węzeł gładko typu slice ma
sygnaturę równą zero.
\end{twierdzenie}

\begin{uwaga}[Zakres dowodu nierówności]
Powyższa nierówność nie wynika z samej
definicji macierzy Seiferta ani z ograniczenia
$|\sigma(K)|\le2g(K)$: jej dowód wykorzystuje
czterowymiarową topologię, np. sygnaturę
dwukrotnego nakrycia $B^4$ rozgałęzionego
wzdłuż właściwie zanurzonej powierzchni
o brzegu $K$ oraz oszacowanie rzędu jego
formy przecięcia. Jest to klasyczny
wynik Murasugiego, przyjmowany tutaj
z literatury. Użycie tej nierówności
w przykładzie trójlistnika nie jest
przedstawiane jako dowód z poprzednich
stwierdzeń o rodzaju trójwymiarowym.
\end{uwaga}

\begin{uwaga}[Sygnatury Levine'a--Tristrama]
Pełniejszy układ niezmienników uzyskuje się z hermitowskich
macierzy $H_\omega:=(1-\omega)V+(1-\overline\omega)V^T$
dla $\omega\in S^1\setminus\{1\}$. Ich sygnatury, gdy macierz
jest niezdegenerowana, definiują funkcję sygnatur węzła.
Dla $\omega=-1$ otrzymujemy
$H_{-1}=2(V+V^T)$, a więc sygnaturę klasyczną.
Przy pierwiastkach wielomianu Alexandera na okręgu jednostkowym
mogą pojawiać się skoki funkcji sygnatur.
\end{uwaga}

\section{Wielomian Jonesa i nawias Kauffmana}
\label{sec:kt-jones}

Wielomian Jonesa jest pierwszym przykładem niezmiennika zbudowanego
wprost z lokalnych reguł rozwiązywania skrzyżowań, który wykracza
poza klasyczną teorię modułu Alexandera. Jego odkrycie wyrosło
z teorii algebr operatorowych, w szczególności z badań podczynników
algebr von Neumanna.

\begin{definicja}[Nawias Kauffmana]
Każdemu niezorientowanemu diagramowi splotu $D$ przyporządkowujemy
wielomian Laurenta $\langle D\rangle\in\mathbb Z[A,A^{-1}]$
rekurencyjnie przez trzy reguły:
\begin{enumerate}[label=(\roman*)]
\item $\langle\bigcirc\rangle=1$ dla diagramu pojedynczego okręgu;
\item $\langle D\sqcup\bigcirc\rangle=
(-A^2-A^{-2})\langle D\rangle$ dla diagramu zawierającego
co najmniej jedną składową;
\item dla lokalnego skrzyżowania
$\langle D_{\times}\rangle
=A\langle D_{\mathrm A}\rangle+A^{-1}\langle D_{\mathrm B}\rangle$,
gdzie $D_{\mathrm A},D_{\mathrm B}$ oznaczają dwa lokalne
wygładzenia w przedstawionej poniżej konwencji. W pokazanym
skrzyżowaniu pasmo biegnące z lewego górnego końca do prawego
dolnego przechodzi górą; wygładzenie A łączy lewy górny
koniec z lewym dolnym oraz prawy górny z prawym dolnym,
a wygładzenie B łączy oba końce górne i oba dolne.
Jeżeli zmienimy, które pasmo przechodzi górą, oznaczenia
A i B zamieniamy miejscami. Reguła nie zależy od orientacji
pasm i jest zgodna z obrotem lokalnego diagramu.
\end{enumerate}
Dla diagramu bez skrzyżowań o $m\geq1$ składowych otrzymujemy
$\langle D\rangle=(-A^2-A^{-2})^{m-1}$.
\end{definicja}

\begin{center}
\begin{tikzpicture}[scale=.86,line cap=round,line join=round,>=Stealth]
\begin{scope}[shift={(-4.0,0)}]
  \draw[ktBlue,line width=2.8pt] (-.73,1)--(.73,-1);
  \draw[ktCoral,line width=2.8pt] (.73,1)--(.13,.18);
  \draw[ktCoral,line width=2.8pt] (-.13,-.18)--(-.73,-1);
  \node[text=ktInk,font=\small] at (0,-1.4) {skrzyżowanie};
\end{scope}
\node[text=ktGold,font=\large\bfseries] at (-2.15,0) {$A$};
\begin{scope}
  \draw[ktTeal,line width=2.8pt] (-.73,-1)--(-.73,1);
  \draw[ktTeal,line width=2.8pt] (.73,-1)--(.73,1);
  \node[text=ktInk,font=\small] at (0,-1.4) {wygładzenie A};
\end{scope}
\node[text=ktGold,font=\large\bfseries] at (2.14,0) {$A^{-1}$};
\begin{scope}[shift={(4.0,0)}]
  \draw[ktPurple,line width=2.8pt] (-.73,1).. controls (-.73,.34) and (.73,.34) .. (.73,1);
  \draw[ktPurple,line width=2.8pt] (-.73,-1).. controls (-.73,-.34) and (.73,-.34) .. (.73,-1);
  \node[text=ktInk,font=\small] at (0,-1.4) {wygładzenie B};
\end{scope}
\end{tikzpicture}
\end{center}

Definicja jest wykonalna: przy $n$ skrzyżowaniach rozwijamy
rekurencję do najwyżej $2^n$ stanów wygładzonych. Dla każdego
stanu $s$ niech $a(s)$ oznacza liczbę wygładzeń typu A,
$b(s)$ liczbę wygładzeń typu B, a $|s|$ liczbę okręgów po
wygładzeniu. Wówczas
\[
\langle D\rangle=\sum_s A^{a(s)-b(s)}
(-A^2-A^{-2})^{|s|-1}.
\]
To pierwszy całkowicie jawny wzór obliczeniowy.

\begin{stwierdzenie}
Nawias Kauffmana jest niezmienny przy ruchach Reidemeistera
$\mathrm{R II}$ i $\mathrm{R III}$. Przy dodatnim lub ujemnym
ruchu $\mathrm{R I}$ zostaje pomnożony odpowiednio przez
$-A^3$ lub $-A^{-3}$, przy zgodnej konwencji znaku skrzyżowania.
\end{stwierdzenie}

\begin{proof}
Przy ruchu $\mathrm{R I}$ jednorazowe rozwinięcie
skrzyżowania i użycie relacji dla dodatkowego okręgu
daje $A(-A^2-A^{-2})+A^{-1}=-A^3$.
Przeciwny skręt daje analogicznie $-A^{-3}$.

Aby dokładnie sprawdzić ruchy $\mathrm{R II}$ i
$\mathrm{R III}$, przejdźmy do lokalnej algebry
diagramów o odpowiednio dwóch i trzech pasmach.
Oznaczmy przez $1$ diagram pionowych pasm,
a przez $e_i$ diagram wygładzenia łączącego
pasma $i$ oraz $i+1$ na górze i na dole.
Składanie diagramów polega na ustawianiu ich
jeden nad drugim; zamknięta pętla wewnętrzna
daje czynnik $\delta=-A^2-A^{-2}$.
Bezpośrednie sklejenie daje relacje
\[
e_i^2=\delta e_i,\qquad
e_1e_2e_1=e_1,\qquad
e_2e_1e_2=e_2.
\]
Lokalne dodatnie i ujemne skrzyżowania
reprezentują odpowiednio
$b_i=A1+A^{-1}e_i$ oraz
$b_i^{-1}=A^{-1}1+Ae_i$.
Dla dwóch pasm obliczamy
\[
b_ib_i^{-1}
=1+(A^2+A^{-2}+\delta)e_i=1.
\]
To jest dokładnie relacja ruchu $\mathrm{R II}$,
włącznie z wszystkimi czterema wygładzeniami.

Dla trzech pasm rozwijamy trzy czynniki i,
korzystając z relacji $e_i^2=\delta e_i$,
$e_1e_2e_1=e_1$, otrzymujemy
\[
\begin{aligned}
b_1b_2b_1
&=A^3 1
+A(2e_1+e_2)
+A^{-1}(e_1e_2+e_2e_1+\delta e_1)
+A^{-3}e_1\\
&=A^3 1+A(e_1+e_2)
+A^{-1}(e_1e_2+e_2e_1).
\end{aligned}
\]
Zamieniając indeksy $1$ i $2$ oraz stosując
$e_2e_1e_2=e_2$, dostajemy dokładnie ten
sam wynik dla $b_2b_1b_2$.
A więc $b_1b_2b_1=b_2b_1b_2$,
co realizuje ruch $\mathrm{R III}$.
Każda z tych równości jest lokalną równością
diagramów z tymi samymi końcami pasm, toteż
pozostaje prawdziwa po dołączeniu dowolnej
pozostałej części diagramu i jego zamknięciu.
Dowiedliśmy niezmienniczości nawiasu przy
obu ruchach bez odwołania do niewypisanych
rachunków na rysunkach.
\end{proof}

\begin{definicja}[Wielomian Jonesa]
Niech $D$ będzie zorientowanym diagramem węzła lub splotu
$L\subset S^3$, a $w(D)$ jego skrętem. Kładziemy
$f_D(A):=(-A^3)^{-w(D)}\langle D\rangle$.
Następnie definiujemy $V_L(t):=f_D(t^{-1/4})$,
rozumiejąc wynik jako wielomian Laurenta w $t^{1/2}$.
Dla węzła przy typowej normalizacji wykładniki są całkowite;
dla splotów mogą występować połówkowe.
\end{definicja}

\begin{twierdzenie}
Wielomian $V_L(t)$ nie zależy od diagramu zorientowanego
splotu $L\subset S^3$. W szczególności $V_U(t)=1$.
Przy konsekwentnie ustalonych konwencjach skrzyżowań spełnia
relację skein
$t^{-1}V_{L_+}(t)-tV_{L_-}(t)=
(t^{1/2}-t^{-1/2})V_{L_0}(t)$,
gdzie $L_+,L_-,L_0$ są jednakowymi poza dyskiem diagramami:
$L_+$ i $L_-$ różnią się znakiem skrzyżowania, zaś $L_0$
powstaje przez zorientowane wygładzenie.
\end{twierdzenie}

\begin{proof}
Wszystkie trzy ruchy Reidemeistera zachowują $f_D$.
Dla $\mathrm{R II}$ i $\mathrm{R III}$ nie zmienia się ani
nawias, ani skręt. Dla $\mathrm{R I}$ mnożnik $-A^{\pm3}$
w nawiasie skraca się z przyrostem $(-A^3)^{-w(D)}$.
Twierdzenie Reidemeistera zapewnia niezależność od diagramu.
Relację skein otrzymujemy, rozwijając lokalnie skrzyżowanie
dla $L_+$ i $L_-$ oraz eliminując oba wygładzenia z
uwzględnieniem różnicy ich skrętów. Po podstawieniu
$t=A^{-4}$ otrzymujemy wskazaną tożsamość.
\end{proof}

\begin{uwaga}[Orientacje i konwencje]
W literaturze spotyka się także zmienną $t^{-1}$ zamiast $t$,
co zamienia wartości dla odbić lustrzanych. Zawsze należy
sprawdzać przyjętą relację skein i normalizację węzła
trywialnego przed porównaniem tablic wielomianów.
\end{uwaga}

\begin{stwierdzenie}
Niech $K\subset S^3$ będzie węzłem, a $\overline K$ jego
odbiciem lustrzanym. Wówczas
$V_{\overline K}(t)=V_K(t^{-1})$.
\end{stwierdzenie}

\begin{proof}
Odbicie lustrzane wymienia wygładzenia A i B, więc
$\langle\overline D\rangle(A)=\langle D\rangle(A^{-1})$.
Jednocześnie $w(\overline D)=-w(D)$.
Po normalizacji i zamianie $A$ na $t^{-1/4}$ dostajemy wynik.
\end{proof}

\begin{przyklad}[Obliczenie wielomianu Jonesa trójlistnika]
Trójlistnik można przedstawić jako domknięcie dodatniego warkocza
$\sigma_1^3$. W dwupasmowej algebrze Temperleya--Lieba oznaczmy przez
$1$ diagram identyczności, a przez $e$ element z jednym połączeniem
między pasmami. Dla generatora warkoczowego mamy wzór
$b=A1+A^{-1}e$, relację $e^2=\delta e$ oraz
$\delta=-A^2-A^{-2}$.
Rozwijając $b^3$, dostajemy
\[
b^3=A^3 1+(A-A^{-3}+A^{-7})e.
\]
Po domknięciu warkocza diagram $1$ daje dwa okręgi, natomiast
$e$ daje jeden okrąg. Z definicji nawiasu Kauffmana wynika zatem
\[
\langle\widehat{\sigma_1^3}\rangle
=A^3\delta+(A-A^{-3}+A^{-7})
=-A^5-A^{-3}+A^{-7}.
\]
Skręt standardowego dodatniego diagramu trójlistnika wynosi $w=3$,
więc po normalizacji otrzymujemy
\[
f(A)=(-A^3)^{-3}\langle\widehat{\sigma_1^3}\rangle
=A^{-4}+A^{-12}-A^{-16}.
\]
Po podstawieniu $t=A^{-4}$ dostajemy
\[
V_T(t)=t+t^3-t^4.
\]
Dla odbicia lustrzanego dostajemy z poprzedniego stwierdzenia
$V_{\overline T}(t)=t^{-1}+t^{-3}-t^{-4}$. Wielomian Jonesa
rozróżnia więc oba odbicia trójlistnika.
\end{przyklad}

\section{Warkocze i grupy warkoczowe}
\label{sec:kt-warkocze}

\begin{definicja}[Warkocz geometryczny]
Niech $n\geq1$. \emph{Warkoczem o $n$ pasmach} nazywamy
zbiór $n$ parami rozłącznych, gładkich łuków w
$\mathbb R^2\times[0,1]$, których początki są ustalonymi
$n$ punktami w $\mathbb R^2\times\{0\}$, końce stanowią
permutację tych punktów w $\mathbb R^2\times\{1\}$,
a rzutowanie każdego łuku na współrzędną $[0,1]$
jest dyfeomorfizmem. Warkocze utożsamiamy z dokładnością
do izotopii zachowującej końce i warunek monotoniczności.
Składanie przez sklejanie końców i przeskalowanie wysokości
tworzy \emph{grupę warkoczową} $B_n$.
\end{definicja}

\begin{twierdzenie}[Prezentacja Artina]
Dla $n\geq2$ grupa warkoczowa $B_n$ ma prezentację
\[
B_n=\Big\langle\sigma_1,\ldots,\sigma_{n-1}\ \Big|\ 
\sigma_i\sigma_j=\sigma_j\sigma_i\ (|i-j|\geq2),\quad
\sigma_i\sigma_{i+1}\sigma_i=\sigma_{i+1}\sigma_i\sigma_{i+1}
\ (1\leq i\leq n-2)\Big\rangle.
\]
Generator $\sigma_i$ oznacza dodatnie skrzyżowanie pasm
o numerach $i$ oraz $i+1$; $\sigma_i^{-1}$ oznacza
skrzyżowanie przeciwne.
\end{twierdzenie}

\begin{proof}[Idea dowodu i wykorzystywane narzędzie]
Każdy warkocz można zaburzyć tak, aby jego rzut na płaszczyznę
poziomej i pionowej współrzędnej rzutowania miał skończenie wiele
skrzyżowań rozdzielonych wysokością. Odczytanie ich w kolejności
daje słowo w generatorach $\sigma_i^{\pm1}$.
W rodzinie izotopii ogólnego położenia zmiany słowa pochodzą
z odwrócenia sąsiednich skrzyżowań, jeśli dotyczą rozłącznych
par pasm, oraz z lokalnego ruchu trzech pasm. Dają one
dokładnie relacje komutacji i relację Artina. Szczegółowe
uzasadnienie zgodności z izotopią względną względem końców
wymaga argumentu o transwersalności.
\end{proof}

\begin{definicja}[Domknięcie warkocza]
Dla $\beta\in B_n$ \emph{domknięciem} $\widehat\beta\subset S^3$
nazywamy splot otrzymany przez połączenie każdego końca
górnego z odpowiadającym mu początkiem dolnym rozłącznymi
łukami biegnącymi poza cylindrem zawierającym warkocz.
\end{definicja}

\begin{twierdzenie}[Alexander]
Każdy zorientowany splot $L\subset S^3$ jest równoważny
domknięciu pewnego warkocza $\beta\in B_n$ dla pewnego $n\geq1$.
\end{twierdzenie}

\begin{proof}[Idea dowodu]
Ustalamy oś w $\mathbb R^3$ rozłączną z $L$ i rozważamy
współrzędną kątową wokół osi. Jeżeli podczas śledzenia splotu
kąt czasem maleje, zastępujemy odpowiednie fragmenty łukami
przebiegającymi po dodatniej stronie osi w taki sposób,
aby stał się lokalnie rosnący. Ta operacja jest wykonywana
skończenie wiele razy i nie zmienia izotopijnego typu splotu.
Po jej zakończeniu współrzędna kątowa jest regularna
na każdej składowej i definiuje zamknięty warkocz.
\end{proof}

\begin{twierdzenie}[Markow]
Niech $\alpha\in B_m$ i $\beta\in B_n$ będą warkoczami.
Ich domknięcia $\widehat\alpha$ i $\widehat\beta$ są
równoważnymi zorientowanymi splotami wtedy i tylko wtedy,
gdy można przejść od jednego warkocza do drugiego skończonym
ciągiem następujących operacji i ich odwrotności:
\begin{enumerate}[label=(\roman*)]
\item sprzęganie $\gamma\mapsto\eta\gamma\eta^{-1}$ wewnątrz $B_k$;
\item stabilizacja $\gamma\in B_k\mapsto\gamma\sigma_k^{\pm1}
\in B_{k+1}$, przy naturalnym włożeniu $B_k\hookrightarrow B_{k+1}$.
\end{enumerate}
\end{twierdzenie}

\begin{uwaga}
Twierdzenia Alexandera i Markowa są podstawą algebraicznych
konstrukcji niezmienników splotów: szukamy reprezentacji
$\rho_n\colon B_n\to\mathrm{GL}(V_n)$ i funkcjonałów typu
śladu, które są zgodne ze sprzęganiem i stabilizacją.
Tak właśnie powstają wielomian Jonesa i pokrewne niezmienniki
w teorii reprezentacji algebr Temperleya--Lieba i Hecke'go.
\end{uwaga}

\begin{center}
\begin{tikzpicture}[line cap=round,scale=.95]
  \begin{scope}[shift={(-2.5,0)}]
    \fill[ktBlue!6,rounded corners=5pt] (-1.35,-.7) rectangle (1.35,2.85);
    \draw[draw=ktBlue,line width=2.4pt] (-0.45,2.45) .. controls (-1.45,2.65) and (-1.45,0.0) .. (-0.45,0.0);
    \draw[draw=ktCoral,line width=2.4pt] ( 0.45,2.45) .. controls ( 1.45,2.65) and ( 1.45,0.0) .. ( 0.45,0.0);
    \foreach \yy in {2.4,1.6,0.8}{
      \draw[draw=ktBlue,line width=2.4pt] (-0.45,\yy)--(0.45,\yy-0.65);
      \draw[draw=ktCoral,line width=2.4pt] ( 0.45,\yy)--(0.12,\yy-0.24);
      \draw[draw=ktCoral,line width=2.4pt] (-0.12,\yy-0.41)--(-0.45,\yy-0.65);
    }
    \foreach \yt/\yb in {2.45/2.4,1.75/1.6,0.95/0.8}{
      \draw[ktTeal,line width=2.4pt] (-0.45,\yb)--(-0.45,\yt);
      \draw[ktGold,line width=2.4pt] ( 0.45,\yb)--( 0.45,\yt);
    }
    \node at (0,-.47) {$\widehat{\sigma_1^3}$: trójlistnik};
  \end{scope}
  \begin{scope}[shift={(2.4,0)}]
    \fill[ktCoral!6,rounded corners=5pt] (-1.35,-.7) rectangle (1.35,2.85);
    \draw[ktPurple,line width=2.4pt] (-0.45,1.65) .. controls (-1.3,1.95) and (-1.3,0.0) .. (-0.45,0.0);
    \draw[ktCoral,line width=2.4pt] ( 0.45,1.65) .. controls ( 1.3,1.95) and ( 1.3,0.0) .. ( 0.45,0.0);
    \foreach \yy in {1.6,0.8}{
      \draw[ktPurple,line width=2.4pt] (-0.45,\yy)--(0.45,\yy-0.65);
      \draw[ktCoral,line width=2.4pt] (0.45,\yy)--(0.12,\yy-0.24);
      \draw[ktCoral,line width=2.4pt] (-0.12,\yy-0.41)--(-0.45,\yy-0.65);
    }
    \foreach \yt/\yb in {1.65/1.6,0.95/0.8}{
      \draw[ktTeal,line width=2.4pt] (-0.45,\yb)--(-0.45,\yt);
      \draw[ktGold,line width=2.4pt] (0.45,\yb)--(0.45,\yt);
    }
    \node at (0,-.47) {$\widehat{\sigma_1^2}$: splot Hopfa};
  \end{scope}
\end{tikzpicture}
\end{center}

\begin{przyklad}[Trójlistnik jako zamknięty warkocz]
Domknięcie warkocza $\sigma_1^3\in B_2\cong\mathbb Z$
jest jednym z dwóch trójlistników. Domknięcie
$\sigma_1^{-3}$ jest jego odbiciem lustrzanym.
Dla $\sigma_1^{2k+1}$ otrzymujemy węzeł torusowy
$T(2,2k+1)$, z orientacją zależną od konwencji znaków.
\end{przyklad}

\section{Suma spójna, węzły pierwsze i rozkład Schuberta}
\label{sec:kt-suma}

\begin{definicja}[Suma spójna]
Niech $K_1,K_2\subset S^3$ będą zorientowanymi węzłami.
Wybieramy dwie rozłączne kule $B_1,B_2$, których przecięcia
z odpowiednimi węzłami są trywialnymi, właściwie zanurzonymi
łukami. Usuwamy wnętrza tych łuków i łączymy ich cztery końce
dwoma rozłącznymi łukami zgodnie z orientacjami. Otrzymany
węzeł oznaczamy $K_1\#K_2$ i nazywamy \emph{sumą spójną}.
Węzeł jest \emph{pierwszy}, jeśli jest nietrywialny i z równości
$K=K_1\#K_2$ wynika, że jeden ze składników jest trywialny.
\end{definicja}

\begin{stwierdzenie}
Suma spójna jest dobrze określona z dokładnością do izotopii,
łączna, przemienna i ma element neutralny $U$.
Ponadto $\Delta_{K_1\#K_2}(t)=\Delta_{K_1}(t)\Delta_{K_2}(t)$
z dokładnością do jednostki $\pm t^k$, a
$\sigma(K_1\#K_2)=\sigma(K_1)+\sigma(K_2)$.
\end{stwierdzenie}

\begin{proof}
Niezależność od punktów połączenia wynika z możliwości
przesuwania krótkich łuków wzdłuż zorientowanych składników.
Łączność i przemienność uzyskujemy przez zmianę położenia
kul łączących. Z powierzchni Seiferta $F_1,F_2$ budujemy
powierzchnię $F_1\mathbin{\natural}F_2$ przez sumę brzegową.
W odpowiedniej bazie jej macierz Seiferta jest blokowo diagonalna:
$V=V_1\oplus V_2$. Otrzymujemy zatem
$\det(tV-V^T)=\det(tV_1-V_1^T)\det(tV_2-V_2^T)$
i analogiczną addytywność sygnatury symetryzacji.
\end{proof}

\begin{twierdzenie}[Schubert; rozkład na czynniki pierwsze]
Każdy nietrywialny węzeł $K\subset S^3$ można przedstawić
jako skończoną sumę spójną węzłów pierwszych. Dwa takie
rozkłady mają tę samą liczbę czynników, a po ich odpowiednim
przestawieniu czynniki odpowiadające sobie są równoważne.
\end{twierdzenie}

\begin{uwaga}
Dowód twierdzenia Schuberta wykorzystuje sfery dekomponujące,
które przecinają węzeł w dwóch punktach, i analizę ich
wzajemnego położenia. Istnienie rozkładu i jego skończoność można uzasadniać
za pomocą teorii sfer dekomponujących i twierdzeń o skończoności
dla powierzchni istotnych w trójrozmaitościach; jednoznaczność
wymaga analizy przecięć takich sfer. Nie jest to
bezpośredni wniosek z jedyności faktoryzacji wielomianów
Alexandera, ponieważ różne węzły mogą mieć ten sam wielomian.
\end{uwaga}

\begin{twierdzenie}[Addytywność rodzaju; Schubert]
Dla dowolnych węzłów $K_1,K_2\subset S^3$ zachodzi
$g(K_1\#K_2)=g(K_1)+g(K_2)$.
\end{twierdzenie}

\begin{proof}[Redukcja do rozcięcia minimalnej powierzchni]
Biorąc powierzchnie Seiferta $F_1,F_2$ minimalnego rodzaju
oraz ich sumę brzegową $F_1\mathbin{\natural}F_2$,
uzyskujemy powierzchnię o jednym brzegu i
$\chi(F_1\mathbin{\natural}F_2)=\chi(F_1)+\chi(F_2)-1$.
Ponieważ $\chi(F_i)=1-2g(F_i)$, otrzymujemy
$g(K_1\#K_2)\le g(K_1)+g(K_2)$.

Niech teraz $F$ będzie powierzchnią Seiferta minimalnego rodzaju
dla $K_1\#K_2$. Wybierzmy sferę dekomponującą
$S\subset S^3$, która spotyka węzeł dokładnie w dwóch
punktach i rozdziela odpowiednie składniki. Ustawiamy $F$
i $S$ poprzecznie. Powierzchnia $F\cap S$ jest sumą
okręgów oraz właściwie zanurzonych łuków z końcami w dwóch
punktach $S\cap\partial F$; jedyny taki łuk musi łączyć
te punkty.

Eliminujemy składowe okręgowe metodą najbardziej wewnętrznego
dysku. Jeśli okrąg $C\subset F\cap S$ ogranicza dysk na
$F$, wybieramy dysk najbardziej wewnętrzny i usuwamy
$C$ przez izotopię. Gdy pozostałby okrąg istotny na $F$,
to rozcinając sferę $S$ wzdłuż najbardziej wewnętrznego
okręgu na $S$ i stosując nieściśliwość powierzchni
minimalnego rodzaju w zewnętrzu węzła, uzyskalibyśmy
sprzeczność. Wykorzystaliśmy tu standardowy lemat:
\emph{powierzchnia Seiferta minimalnego rodzaju jest
nieściśliwa w zewnętrzu węzła}. Jego dowód polega na
kompresji: istotny dysk kompresujący zmniejsza rodzaj,
nie zmieniając zadanej składowej brzegu.

Po usunięciu okręgów $F\cap S$ pozostaje jeden łuk.
Rozcinamy $F$ wzdłuż tego łuku, a brzegi powstałe na
sferze zaślepiamy odpowiednimi dyskami po obu stronach
$S$. Dostajemy dwie powierzchnie Seiferta $F_1',F_2'$
dla $K_1,K_2$. Rozcięcie wzdłuż łuku zwiększa
charakterystykę Eulera o jeden:
$\chi(F_1')+\chi(F_2')=\chi(F)+1$.
Ponieważ wszystkie trzy powierzchnie są spójne i mają
po jednej składowej brzegu, daje to
$g(F)=g(F_1')+g(F_2')\ge g(K_1)+g(K_2)$.
W połączeniu z poprzednią nierównością kończy to dowód.
\end{proof}

\section{Węzły torusowe, włóknistość i rodzaj}
\label{sec:kt-torusowe}

\begin{twierdzenie}[Grupa węzła torusowego]
Niech $p,q\geq2$ będą względnie pierwsze. Dla węzła torusowego
$T(p,q)\subset S^3$ zachodzi
$\pi_1(S^3\setminus T(p,q))\cong
\langle a,b\mid a^p=b^q\rangle$.
\end{twierdzenie}

\begin{proof}
Standardowy torus rozdziela $S^3$ na dwa torusy pełne
$V_1,V_2$. Węzeł $T(p,q)$ jest istotną krzywą na ich wspólnym
brzegu. Po wycięciu wąskiego otoczenia tej krzywej
uzupełnienie węzła jest sumą dwóch torusów pełnych sklejonych
wzdłuż pierścienia. Rdzenie torusów pełnych generują
$\pi_1(V_1)\cong\mathbb Z$ i $\pi_1(V_2)\cong\mathbb Z$.
Pierścień sklejający biegnie odpowiednio $p$ i $q$ razy
wokół tych rdzeni, zatem twierdzenie van Kampena daje
prezentację z relacją $a^p=b^q$.
\end{proof}

\begin{twierdzenie}[Rodzaj węzła torusowego]
Jeżeli $p,q\geq2$ są względnie pierwsze, to
$g(T(p,q))=(p-1)(q-1)/2$.
\end{twierdzenie}

\begin{proof}
Dodatni warkocz $\beta=(\sigma_1\dots\sigma_{p-1})^q$
ma domknięcie równe $T(p,q)$. Ma $(p-1)q$ skrzyżowań
i $p$ okręgów Seiferta, więc algorytm Seiferta daje
powierzchnię rodzaju $(1-p+(p-1)q)/2=(p-1)(q-1)/2$.
Dla oszacowania z dołu obliczamy z prezentacji grupy
węzła (lub z odpowiedniej macierzy Seiferta) wielomian
\[
\Delta_{T(p,q)}(t)=
\frac{(t^{pq}-1)(t-1)}{(t^p-1)(t^q-1)}
\]
z dokładnością do jednostki $\pm t^k$.
Licznik i mianownik skracają się do wielomianu o szerokości
wykładników $(p-1)(q-1)$, ponieważ po redukcji stopień
wynosi $pq+1-p-q$. Z nierówności
$\operatorname{span}\Delta_K\leq2g(K)$ otrzymujemy
$g(T(p,q))\geq(p-1)(q-1)/2$.
\end{proof}

\begin{definicja}[Węzeł włóknisty]
Węzeł $K\subset S^3$ nazywamy \emph{włóknistym}, jeżeli jego
uzupełnienie $X_K=S^3\setminus\operatorname{int}N(K)$
jest przestrzenią wiązki nad $S^1$ o włóknie będącym zwartą,
spójną powierzchnią $F$ z jednym brzegiem, przy czym
$\partial F$ odpowiada preferowanej długości na
$\partial X_K$. Dla pewnego dyfeomorfizmu
$\varphi\colon F\to F$ identycznego w pobliżu brzegu mamy
$X_K\cong (F\times[0,1])/(x,1)\sim(\varphi(x),0)$.
Dyfeomorfizm $\varphi$ nazywamy \emph{monodromią}.
\end{definicja}

\begin{twierdzenie}
Jeżeli $K\subset S^3$ jest węzłem włóknistym, to jego
wielomian Alexandera jest jednomianowy na obu krańcach,
tzn. współczynniki przy najwyższej i najniższej potędze
$t$ są równe $\pm1$. Każdy węzeł torusowy jest włóknisty.
\end{twierdzenie}

\begin{proof}[Uzasadnienie pierwszej tezy]
Dla węzła włóknistego nakrycie nieskończenie cykliczne
uzupełnienia jest homotopijnie równoważne włóknu $F$.
Na $H_1(F;\mathbb Z)$ działanie generatora przekształceń
pokładowych jest dane przez automorfizm całkowity
$\varphi_*\in GL_{2g}(\mathbb Z)$. Wielomian Alexandera
jest, z dokładnością do jednostki, wielomianem
charakterystycznym tego automorfizmu:
$\Delta_K(t)\doteq\det(tI-\varphi_*)$.
Jego najwyższy współczynnik wynosi jeden, a wyraz wolny
$\det(-\varphi_*)=\pm1$. Włóknistość węzłów torusowych
wynika z opisu ich uzupełnień jako rozwarstwienia z monodromią;
pełny dowód można otrzymać przez analizę osobliwości
$f(z,w)=z^p+w^q$ i fibration Milnora na małej sferze
w $\mathbb C^2$.
\end{proof}

\begin{uwaga}
Jednomianowość wielomianu Alexandera jest warunkiem koniecznym,
ale nie wystarczającym włóknistości. Później zobaczymy, że
homologia Floera węzła pozwala rozpoznawać włóknistość,
a nie tylko wykrywać tę przeszkodę.
\end{uwaga}



